% concatenated from main.tex
%%%%%%%% ICML 2025 EXAMPLE LATEX SUBMISSION FILE %%%%%%%%%%%%%%%%%

\documentclass{article}
% Recommended, but optional, packages for figures and better typesetting:
\usepackage{microtype}
\usepackage{graphicx}
\usepackage{subfigure}
\usepackage{booktabs} % for professional tables
\usepackage{amsmath}

\usepackage{floatflt} 
\usepackage{adjustbox}
\usepackage{float}

\newcommand{\cA}{\mathcal{A}}
\newcommand{\cB}{\mathcal{B}}
\newcommand{\cC}{\mathcal{C}}
\newcommand{\cD}{\mathcal{D}}
\newcommand{\cE}{\mathbb{E}}
\newcommand{\cF}{\mathcal{F}}
\newcommand{\cG}{\mathcal{G}}
\newcommand{\cH}{\mathcal{H}}
\newcommand{\cI}{\mathcal{I}}
\newcommand{\cJ}{\mathcal{J}}
\newcommand{\cK}{\mathcal{K}}
\newcommand{\cL}{\mathcal{L}}
\newcommand{\cM}{\mathcal{M}}
\newcommand{\cN}{\mathcal{N}}
\newcommand{\cO}{\mathcal{O}}
\newcommand{\cP}{\mathbb{P}}
\newcommand{\cQ}{\mathcal{Q}}
\newcommand{\cR}{\mathcal{R}}
\newcommand{\cS}{\mathcal{S}}
\newcommand{\cT}{\mathcal{T}}
\newcommand{\cU}{\mathcal{U}}
\newcommand{\cV}{\mathcal{V}}
\newcommand{\cW}{\mathcal{W}}
\newcommand{\cX}{\mathcal{X}}
\newcommand{\cY}{\mathcal{Y}}
\newcommand{\cZ}{\mathcal{Z}}
%% math blackboard bold
\newcommand{\A}{\mathbb{A}}
\newcommand{\B}{\mathbb{B}}
\newcommand{\C}{\mathbb{C}}
\newcommand{\D}{\mathbb{D}}
\newcommand{\E}{\mathbb{E}}
\newcommand{\F}{\mathbb{F}}
\newcommand{\G}{\mathbb{G}}
\newcommand{\I}{\mathbb{I}}
\newcommand{\J}{\mathbb{J}}
\newcommand{\K}{\mathbb{K}}
\newcommand{\M}{\mathbb{M}}
\newcommand{\N}{\mathbb{N}}
\renewcommand{\Pr}{\mathbb{P}}
\renewcommand{\S}{\mathbb{S}}
\newcommand{\Q}{\mathbb{Q}}
\newcommand{\R}{\mathbb{R}}
\newcommand{\T}{\mathbb{T}}
\newcommand{\U}{\mathbb{U}}
\newcommand{\V}{\mathbb{V}}
\newcommand{\W}{\mathbb{W}}
\newcommand{\X}{\mathbb{X}}
\newcommand{\Y}{\mathbb{Y}}
\newcommand{\Z}{\mathbb{Z}}

\DeclareMathOperator{\argmax}{arg\,max}
\DeclareMathOperator{\argmin}{arg\,min}

% hyperref makes hyperlinks in the resulting PDF.
% If your build breaks (sometimes temporarily if a hyperlink spans a page)
% please comment out the following usepackage line and replace
% \usepackage{icml2025} with \usepackage[nohyperref]{icml2025} above.
\usepackage{hyperref}


% Attempt to make hyperref and algorithmic work together better:
\newcommand{\theHalgorithm}{\arabic{algorithm}}

% Use the following line for the initial blind version submitted for review:
% \usepackage{icml2025}

% If accepted, instead use the following line for the camera-ready submission:
%\usepackage{icml2025}
\usepackage[accepted]{icml2025}

% For theorems and such
\usepackage{amsmath}
\usepackage{amssymb}
\usepackage{mathtools}
\usepackage{amsthm}

% if you use cleveref..
\usepackage[capitalize,noabbrev]{cleveref}

% for orcid-id links
\usepackage{orcidlink}

%%%%%%%%%%%%%%%%%%%%%%%%%%%%%%%%
% THEOREMS
%%%%%%%%%%%%%%%%%%%%%%%%%%%%%%%%
\theoremstyle{plain}
\newtheorem{theorem}{Theorem}[section]
\newtheorem{proposition}[theorem]{Proposition}
\newtheorem{lemma}[theorem]{Lemma}
\newtheorem{corollary}[theorem]{Corollary}
\theoremstyle{definition}
\newtheorem{definition}[theorem]{Definition}
\newtheorem{assumption}[theorem]{Assumption}
\theoremstyle{remark}
\newtheorem{remark}[theorem]{Remark}

% Todonotes is useful during development; simply uncomment the next line
%    and comment out the line below the next line to turn off comments
%\usepackage[disable,textsize=tiny]{todonotes}
\usepackage[textsize=tiny]{todonotes}


% The \icmltitle you define below is probably too long as a header.
% Therefore, a short form for the running title is supplied here:
\icmltitlerunning{ADDQ: Adaptive Distributional Double $Q$-Learning}

\begin{document}

\twocolumn[
\icmltitle{ADDQ: Adaptive Distributional Double $Q$-Learning}

% It is OKAY to include author information, even for blind
% submissions: the style file will automatically remove it for you
% unless you've provided the [accepted] option to the icml2025
% package.

% List of affiliations: The first argument should be a (short)
% identifier you will use later to specify author affiliations
% Academic affiliations should list Department, University, City, Region, Country
% Industry affiliations should list Company, City, Region, Country

% You can specify symbols, otherwise they are numbered in order.
% Ideally, you should not use this facility. Affiliations will be numbered
% in order of appearance and this is the preferred way.
\icmlsetsymbol{equal}{*}

\begin{icmlauthorlist}
\icmlauthor{Leif D\"oring \orcidlink{0000-0002-4569-5083}}{yyy}
\icmlauthor{Benedikt Wille}{yyy}
\icmlauthor{Maximilian Birr}{yyy}
\icmlauthor{Mihail B\^{\i}rsan}{sch}
\icmlauthor{Martin Slowik \orcidlink{0000-0001-5373-5754}}{yyy}
%\icmlauthor{}{sch}
%\icmlauthor{}{sch}
%\icmlauthor{}{sch}
\end{icmlauthorlist}

\icmlaffiliation{yyy}{Institute of Mathematics, University of Mannheim, Germany}
\icmlaffiliation{sch}{Department of Mathematics and Computer Science, Freie Universität Berlin, Germany}

\icmlcorrespondingauthor{Leif D\"oring}{doering@uni-mannheim.de}

% You may provide any keywords that you
% find helpful for describing your paper; these are used to populate
% the "keywords" metadata in the PDF but will not be shown in the document
\icmlkeywords{Machine Learning, ICML}

\vskip 0.3in
]

% this must go after the closing bracket ] following \twocolumn[ ...

% This comm actually creates the footnote in the first column
% listing the affiliations and the copyright notice.
% The command takes one argument, which is text to display at the start of the footnote.
% The \icmlEqualContribution command is standard text for equal contribution.
% Remove it (just {}) if you do not need this facility.

\printAffiliationsAndNotice{}  % leave blank if no need to mention equal contribution
%\printAffiliationsAndNotice{\icmlEqualContribution} % otherwise use the standard text.




\begin{abstract}
  Bias problems in the estimation of $Q$-values are a well-known obstacle that slows down convergence of $Q$-learning and actor-critic methods. One of the reasons of the success of modern RL algorithms is partially a direct or indirect overestimation reduction mechanism.   We propose an easy to implement method built on top of distributional reinforcement learning (DRL) algorithms to deal with the overestimation in a locally adaptive way. Our framework is simple to implement, existing distributional algorithms can be improved with a few lines of code. We provide theoretical evidence and use double $Q$-learning to show how to include locally adaptive overestimation control in existing algorithms. Experiments are provided for tabular, Atari, and MuJoCo environments.
\end{abstract}


\section{Introduction}



A fundamental building block of many modern reinforcement learning (RL) algorithms is Watkins' $Q$-learning (QL) \cite{Watkins1992}. In each round the agent observes a new reward signal and updates the currently estimated state-action function by combining the new reward signal with the best currently estimated action in the next step. Unfortunately, the update rule involves a maximum and maxima suffer both from overestimation bias and function approximation uncertainty. Thus, estimated $Q$ values are initially way too large. Although convergence of $Q$-learning and its variants for tabular cases can be proved rigorously, the convergence can often be seen only after millions of iterations. In the context of $Q$-learning we refer to the seminal paper \cite{Thrun+Schwartz:1993}. The overestimation effect is harmful not only for simple QL and variants with function approximation such as DQN \cite{dqn}, but also for critic estimation in actor-critic methods such as soft actor-critic (SAC) of \cite{sac}. Motivated by statistical approaches to the estimation of the expectation of maxima of random variables the concept of double $Q$-learning (DQL) was introduced in \cite{hasselt2010}. Instead of using one set of random variables two independent sets are used. One is used to detect the maximal index, the other to evaluate the random variable corresponding to the maximal index. For $Q$-learning this translates to keeping track of two copies of the $Q$-matrix that are alternated in order to detect the best action and evaluate the corresponding $Q$-value. DQL and its actor-critic variants reduce the overestimation (see for instance Figure 2 in \cite{fujimoto2018addressing}) and sometimes even underestimate $Q$-values. This, for example, can be seen in a simple chain MDP (Example 6.7 in \cite{Sutton1998} and also \cite{Maxmin} for the underestimation effect in the same example). \cite{fujimoto2018addressing} argue that overestimation should be addressed particularly in state-action regions with high uncertainty. Indirectly, this is taken into account by ensemble methods such as Maxmin QL \cite{Maxmin}. These methods use ensembles of more than two $Q$-estimators. Different ensemble estimators take the minimum, full ensemble averages \cite{averagedQ}, or random ensemble averages \cite{redq}. A related line of research uses uncertainty-based RL, also with the goal of reducing overestimation, see for instance \cite{Wu}, \cite{Ghas}.

There is no rule whether QL, DQL, or any ensemble variant works well for a given environment. Algorithms sometimes perform well and sometimes fail. This article proposes a novel approach. We propose to take QL and (at least) one other method and try to control if the QL update should be replaced or mixed with another underestimating method. The control must be locally adaptive, and the need to manage the bias depends on the local uncertainty (aleatoric and epistemic randomness, function approximation). For an approach similar in spirit to ours, see \cite{Dorka}.
\begin{itemize}
    \item We show theoretically how distributional RL helps the agent identify the need for overestimation control.
    \item As a test case, we combine QL and DQL using a local weighting that we call ADDQ.
    \item Convergence of ADDQ is proved, experiments are performed on tabular, Atari, and MuJoCo environments.
\end{itemize}



\section{$Q$-learning and the overestimation problem}
\subsection{Tabular $Q$-learning}
Let us fix a (discrete) Markov decision model $(\mathcal S,\mathcal A,\mathcal R,p)$, where $\mathcal S$ is a finite state-space, $\mathcal A$ a finite space of allowed actions, $\mathcal R$ the reward space, and $p$ a transition kernel describing the distribution of the reward $r$ and the new state $s'$ when action $a$ is played in state $s$. Given a time-stationary policy $\pi$, a Markov kernel on $\mathcal S\times \mathcal A$, there is a Markov reward process $(S_t,A_t,R_t)$ with transitions 
\begin{align*}
    &\quad \mathbb P^\pi(R_t=r, S_{t+1}=s', A_{t+1}=a'|S_t=s, A_t=a)\\
&=\pi(a'\colon s')p(r,s'\colon s,a).
\end{align*}
The goal of the agent in reinforcement learning is to use  rollouts of the MDP to find a policy that maximizes $Q^\pi(s,a)=\mathbb E^\pi[\sum_{t=0}^\infty \gamma^t R_t|S_0=s,A_0=a]$, the expected discounted reward. The discounting factor $\gamma\in (0,1)$ is fixed. In the discrete setting with $\mathcal S$ and $\mathcal A$ finite it is well-known that optimal stationary policies exist and can be found as greedy policy obtained by the unique solution matrix $Q^*$ to $T^*Q=Q$. The non-linear operator $(T^* Q)(s,a)= r(s,a)+\sum_{s'\in \mathcal S}p(\mathcal R\times \{s'\}:s,a) \gamma\max_{a'\in \mathcal A} Q(s,a')$ is called Bellman's optimality operator. Bellman's optimality operator is a max-norm contraction on the $\mathcal S\times \mathcal A$ matrices. Using Banach's fixed point theorem, the solution can in principle be found by iteratively applying $T^*$ to some initial matrix $Q_0$. The drawback of this approach is the need to know the operator $T^*$, thus having knowledge about the transitions $p$. Using standard stochastic approximation algorithms the fixed point $Q^*$ can be approximated by 
\begin{align*}
     Q(s,a)
    \leftarrow (1-\alpha)Q(s,a)+\alpha\big(r+\gamma \max_{a'}Q(s',a')\big),
\end{align*}
called $Q$-learning (QL). The state-action pairs can be chosen synchronously or asynchronously using rollouts. Typically, to update at $(s,a$) a one-step sample $s',r$ is obtained from $p(\cdot: s,a)$ and the step-sizes $\alpha$ are assumed to satisfy the Robbins-Monro conditions. The exploration (choices of $(s,a)$ to be updated) can be on-policy (using the $Q$-estimates) or off-policy (using a behavior policy). The only requirement is infinite visits for all state-action pairs. The recursively defined matrix-sequence $(Q_t)$ was proved to converge in the tabular setting, see e.g. \cite{Tsitsiklis1994}.
% The problem of $Q$-learning is the tendency for overestimation, the values $Q_t(s,a)$ will typically be (much) larger than $Q_*(s,a)$. We will discuss the rationale behind overestimation and different patches in the sections below. While in statistical procedures a bias is generally unwished for there are situations in which overestimation in $Q$-learning can help. One of several possible reasons is the combination with on-policy exploration where the overestimation can possibly trigger useful exploration. 


\subsection{The overestimation problem}\label{sec:overestimationproblem}
Even though QL converges to $Q^*$ for the number of updates going to infinity, the convergence is very slow. One of the known sources is the so-called overestimation problem of QL. The algorithm does not provide unbiased estimates of $Q^*(s,a)$. Instead, the estimates $Q_t(s,a)$ tend to overestimate $Q^*(s,a)$. A statistical explanation is based on the simple fact that the point estimator $\max\{\hat X_1,...,\hat X_n\}$ is not an unbiased estimator of $\max\{\mathbb E[X_1],...,\mathbb E[X_n]\}$ but the estimator is positively biased. Thus, the update targets $r+\gamma \max_{a'} Q_{t}(s',a')$ can be seen as overestimating the true Bellman optimality operator at each step. The consequences of overestimation are less obvious than it seems at first sight. Shifting $Q$-values the same amount globally has no effect, neither for $Q$-based exploration purposes nor for best action selection. It is the local difference in overestimation that must be avoided to not confuse the agent. Thus, it is crucial to understand the root causes of overestimation to then mitigate by taking alternative updates to QL if needed.

There is little quantitative understanding of the overestimation; for some rough bounds see \cite{VanHasselt:2011}. We used tools from probability theory to compute estimation bounds for a simple explanatory example. 
%
%provide a lower bound for an example (actually a two-sided %version) that was discussed also 
%in Example 6.7 of \cite{Sutton1998} or \cite{Maxmin}.
%
\begin{figure}[h]
\begin{center}
\begin{tikzpicture}[scale=0.8]
 %sutton
    \node[circle, draw] (A) at (6,0) {$s_0$};
    \node[text width=1cm] at (7,0.5) {$r=0$};
    %\node[text width=0.5cm] at (6.5,-0.6) {$r=0$};
    \node[text width=1cm] at (5.3,0.5) {$r=0$};
    \node[text width=2cm] at (3.5,-1.2) {$k_1$ $\mathcal N(\mu_1,\sigma_1^2)$  actions};
    %\node[text width=2cm] at (8.5,-1) {$K_C$ identical};
    \node[circle, draw] (B) at (4,0) {$s_1$};
    \draw[->] (A) -- (B);
    %\node[circle, draw] (C) at (8,0) {C};
% (A) -- (6,-1);
    \filldraw[fill=gray, draw=black] (1.7,-0.25) rectangle +(0.5,0.5);
    %\draw[->] (B) to[out=-105, in=-75]  (2.25,0);
    %\draw[->] (B) to[out=105, in=75]  (2.25,0);
    \draw[->] (B) to[out=135, in=45]  (2.25,0);
    \draw[->] (B) to[out=-135, in=-45]  (2.25,0);
    \fill[black] (3.2,0.25) circle (0.05);
    \fill[black] (3.2,0) circle (0.05);
    \fill[black] (3.2,-0.25) circle (0.05);
    \node[circle, draw] (C) at (8,0) {$s_2$};
    \draw[->] (C) to[out=45, in=135]  (9.65,0);
    \draw[->] (C) to[out=-45, in=-135]  (9.65,0);
    \draw[->] (A) -- (C);

    \filldraw[fill=gray, draw=black] (9.7,-0.25) rectangle +(0.5,0.5);
  \fill[black] (8.7,0.25) circle (0.05);
    \fill[black] (8.7,0) circle (0.05);
    \fill[black] (8.7,-0.25) circle (0.05);
        \node[text width=2cm] at (9,-1.2) {$k_2$ $\mathcal N(\mu_2,\sigma_2^2)$ actions};
    \node[circle, draw] (D) at (6.05,-1.39) {};
    \filldraw[fill=gray, draw=black] (5.8,-1.7) rectangle +(0.5,0.5);
    \draw[->] (A) -- (D);
    \node[text width=1cm] at (6.8,-0.8) {$r=0$};


        
%    \filldraw[fill=gray, draw=black] (7.6,-0.25) rectangle +(0.5,0.5);
    %\draw[->] (C) to[out=-75, in=-105]  (9.6125,0);
    %\draw[->] (C) to[out=75, in=105]  (9.6125,0);
%    \draw[->] (C) to[out=45, in=135]  (9.6125,0);
 %   \draw[->] (C) to[out=-45, in=-135]  (9.6125,0);
  %  \fill[black] (8.8,0.25) circle (0.05);
   % \fill[black] (8.8,0) circle (0.05);
    %\fill[black] (8.8,-0.25) circle (0.05);
\end{tikzpicture}\label{figure71}
\end{center}
\caption{Two-sided bandit MDP, start in $s_0$, gray boxes  terminal}
\end{figure}
The example is intriguing, as it is simple but hard to learn - even more so if one side is replaced by a chain of decisions. Each side gives a reward (for simplicity $0$) followed by a Gaussian reward from one of $k$ actions. QL and DQL can both fail badly (each on one side) for the same parameter configuration. If $\mu_1>0>\mu_2$, then the optimal action in $s_0$ is "left". If $\sigma_2$ (and/or $k_2$) is large compared to $\sigma_1$ (and/or $k_1$) then the overestimated $Q$-values of QL will confuse the agent and lead him to believe that "right" is optimal. Similarly, underestimating can make the agent believe "down" is optimal. Our approach of learning locally to use QL or DQL (or another variant) can mitigate that problem by learning to use QL for one side and DQL for the other. The method is motivated by two theoretical results, a lower bound on overestimation (Proposition \ref{prop2}) and a computation with DRL to connect estimated sample variances and overestimation (Proposition \ref{prop3}). 
%\begin{figure}[h]\label{figure:bandit}
%\begin{center}
%\begin{tikzpicture}[scale=0.8]
 %sutton
 %   \node[circle, draw] (A) at (6,0) {$s_0$};
  %   \node[circle, draw] (C) at (7.8,0) {};
%
 %   \node[text width=0.5cm] at (7,0.5) {0};
    %\node[text width=0.5cm] at (6.5,-0.6) {0};
  %  \node[text width=0.5cm] at (5.3,0.5) {0};
 %   \node[text width=3cm] at (3.5,-1) {$k$ $\mathcal N(\mu,\sigma^2)$ actions};
  %  %\node[text width=2cm] at (8.5,-1) {$K_C$ identical};
 %   \node[circle, draw] (B) at (4,0) {$s_1$};
  %  \draw[->] (A) -- (B);
    %\node[circle, draw] (C) at (8,0) {C};
   %     \draw[->] (A) -- (C);
    %\draw[->] (A) -- (C);
   % \filldraw[fill=gray, draw=black] (1.7,-0.25) rectangle +(0.5,0.5);
   % \filldraw[fill=gray, draw=black] (7.6,-0.25) rectangle +(0.5,0.5);
   %  \draw[->] (B) to[out=135, in=45]  (2.25,0);
   % \draw[->] (B) to[out=-135, in=-45]  (2.25,0);
   % \fill[black] (3.2,0.25) circle (0.05);
   % \fill[black] (3.2,0) circle (0.05);
   % \fill[black] (3.2,-0.25) circle (0.05);
%\end{tikzpicture}
%\caption{Bandit MDP, start in $s_0$, gray boxes are terminating states}
%\end{center}
%\end{figure}
If step-sizes are chosen in the common way as $\alpha_t(s,a)=\frac{1}{T_{s,a}(t)}$, with $T_{s,a}(t)$ the number of visits at $(s,a)$ up to time $t$, then results on sums and maxima of Gaussian random variables can be used to prove a lower bound on the expected overestimation of the true value $\gamma\mu$.
\begin{proposition}\label{prop2}
If the left side has been explored $Nk_1$ times and the exploration was sufficiently exploratory (see Theorem \ref{thm:1}), then the $Q$-estimate at $(s_0,\text{"left"})$ has bias at least $\frac{\gamma}{\sqrt{\pi\log(2) }}\frac{\sigma_1 \sqrt{\log(k_1)}}{\sqrt{N}}$ and analogously at $(s_0,\text{"right"})$.
\end{proposition}
 The lower bound quantifies the idea that uncertainty forces overestimation and the bias only decreases slowly over time. A proof is given in Appendix \ref{appendix:theory}. Since the estimate is relatively tight one gets a feeling for how many reward samples are needed to get sufficiently precise estimates of the $Q$-values so the agent makes the right decision.

There has been considerable interest for the past years to understand the sources of uncertainty in RL. While many sources of uncertainty exist, they are often categorized as either aleatoric or epistemic. Aleatoric uncertainty is model given and cannot be improved by more data or learning. In RL this is uncertainty implied by random variables governing rewards and transitions. Epistemic uncertainty refers to the uncertainty that could potentially be reduced using more data and better algorithms (including better function approximation). Epistemic uncertainty in RL is induced by all random variables to run the learning procedure (exploration, replay buffer, etc.) and function approximation in deep RL. Keeping in mind the different sources of uncertainty is useful in order to identify algorithmic potential for improvement but also theoretical limitations. It is also important to realize that aleatoric and epistemic randomness strongly influence each other. If a reward has large variance (aleatoric uncertainty) then  the estimation with samples creates more epistemic uncertainty. For estimating expectations, this is due to the central limit theorem. 

In deep RL one of the major sources of epistemic uncertainty is function approximation. As in \cite{Thrun+Schwartz:1993} we could add to our analysis independent error-noise modeling the function approximation. If the error noise is assumed Gaussian then our results readily extend, by adding additional noise-variance to our results. Since the modeling assumption as independent noise is rather special we restrain from studying the effect of function approximation on epistemic uncertainty. Instead, we will now show how to use additional information from distributional QL to deal with overestimation in a local way.

\subsection{Tabular distributional $Q$-learning}
To make this article as self-contained as possible we give a minimal  overview on DRL. For a concise treatment we refer to the recent book \cite{drl}. Given a Markov decision model and a stationary policy $\pi$, \cite{rowland18} define the return distribution function as
\[ \eta^{\pi}(s,a)(B):=\mathbb P^{\pi}\Big(\sum_{t=0}^{\infty}\gamma^{t}R_{t}\in B \Big|S_0=s, A_0=a\Big)\]
for $B\in \mathcal B(\mathbb R)$. The expectation over the measure $\eta^\pi$ is the classical state-action value function $Q^\pi$. In contrast to ordinary RL, the target in DRL  is to learn return distributions instead of only return expectations. There have been a lot of theoretical articles on DRL \cite{bellemare2017,dabney2018,rowland18,lyle19,drl,rowland2023,rowland2023analysis} establishing distributional Bellman operators, contractivity, convergence proofs of dynamic programming, etc. It was shown in \cite{bellemare2017,drl} that the return distribution function is the unique solution to $\eta^\pi=T^\pi\eta^\pi,$ where $T^\pi:\mathcal P(\mathbb R)^{\mathcal S\times\mathcal A}\to\mathcal P(\mathbb R)^{\mathcal S\times\mathcal A}$ is the distributional Bellman operator defined as $(T^\pi\eta)(s,a)=\sum_{r,s',a'\in\mathcal R\times\mathcal S\times\mathcal A}b_{r,\gamma}\#\eta(s',a')\pi(a';s')p(s',r;s,a)$ with bootstrap function $b_{r,\gamma}(z)= r+\gamma z$ and  push-forward of measures $f\#\nu(B):=\nu(f^{-1}(B))$. In essence, sample based dynamic programming can also be carried out in a distributional sense replacing the classical Bellman operator by its analogue in the distributional sense. Distributional QL proceeds similarly to classical expectation QL:
\begin{align*}
    \eta(s,a)\leftarrow (1-\alpha) \eta(s,a)+\alpha\big(b_{r,\gamma}\#\eta(s',a^*)\big),
\end{align*}
with $a^*=\operatorname{argmax}_{a'}Q(s',a')$, where $Q(s',a')$ are the expectations of the probability measures $\eta(s',a')$. In order to work algorithmically with DRL parametrizations $\mathcal F$ of measures need to be used. Distributional learning algorithms in practice then work similarly to deep learning algorithms, alternating Bellman operators and function class projection, see Algorithm \ref{alg:CDRL}. 
%%% TODO: eventuell function class projection einführen, comment reviewer 3
\begin{algorithm}[h]
\caption{Distributional $Q$-learning update step}
\label{alg:CDRL}
\begin{algorithmic}
\REQUIRE Proxy $\eta$ for $\eta^*$ and pair $(s,a)$ to be updated
\STATE Determine step-size $\alpha$
\STATE Sample reward/next state $(r, s')$
\STATE {\color{gray} \# Compute target}
\STATE $a^* \gets \operatorname{argmax}_a \mathbb{E}_{Z \sim \eta(s',a)} [Z]$
\STATE $\hat{\eta}_* \gets b_{r, \gamma} \#\eta(s',a^*)$\\
\STATE {\color{gray} \# Project target back onto support}
\STATE $\hat{\eta} \gets \Pi_{\mathcal F} (\hat{\eta}_*)$
\STATE {\color{gray} \# Move $\eta(s,a)$ towards the target, for tabular RL e.g.}
\STATE $\eta(s,a) \leftarrow (1 - \alpha)\eta(s,a) + \alpha\hat{\eta}$
%\STATE $\eta_{new}(\bar s,\bar a)=\eta_{old}(\bar s,\bar a)$ for all $(\bar s,\bar a)\neq(s,a)$\\
%\OUTPUT $\eta^A_{t+1}$
\end{algorithmic}
\end{algorithm}
 There are two simple parametrization that have been used frequently. The categorical parametrization (fixing number of atoms with variable weights at fixed locations) and the quantile parametrization (fixing number of atoms with fixed weights but variable locations). For the categorical parametrization a set of \( m \) evenly spaced locations \( \theta_1 < \cdots < \theta_m \) needs to be fixed, the categorical measures are then defined by $\mathcal{F}_{C,m} = \big\{ \sum_{i=1}^{m} p_i\delta_{\theta_i}\, \big|\, p_i \geq 0, \sum_{i=1}^{m} p_i = 1 \big\}.$ That is, measures in $\mathcal F_{C,m}$ are parametrized by an $m$dim probability vector with weights for the $m$ fixed atoms. In contrast, the quantile parametrization $\mathcal{F}_{Q,m} = \Big\{  \sum_{i=1}^{m}\frac{1}{m} \delta_{\theta_i} : \theta_i \in \mathbb{R} \Big\}$ fixes the weights to $\frac 1 m$ with variable atom locations.

 \subsection{Overestimation mitigation using distributional RL}\label{sec:overestimation}
We follow an insight from Proposition \ref{prop2} that large uncertainty, aleatoric (large $\sigma$) as well as epistemic (small $N$ - additional $\sigma$ if function approximation is modeled as Gaussian error), implies large overestimation of $Q$-values. While we skip function approximation for theoretical considerations we are as precise as possible in the simplest tabular settings. The estimate from Proposition \ref{prop2} suggests to replace QL-updates in $(s,a)$ if uncertainty is large (compared to other actions). Unfortunately, in standard QL the agent has no direct access to such information in order to adjust the update rule at $(s,a)$. This is where our idea comes into play. DRL gives the agent exactly the needed information using distribution insight into the return estimate $\eta_t(s,a)$. It is crucial to note that DRL learns random measures as $\eta(s,a)$ is a probability measure that depends on the random samples used in the updates so it has an expectation and a variance that are both random in terms of the random samples. The situation becomes tricky as we are going to take variances of the variances. To avoid confusion an analogy to statistics is used. We speak of sample averages $M$ (resp. sample variances $S^2$) of $\eta(s,a)$ and expectation $\E$ (resp. variance $\mathbb V$) for the integrals against the randomness induced by the probability space behind all random variables.

We now use the bandit MDP from above to explain why the DRL agent has access to uncertainty estimates during the learning process. To allow concrete computations with distributions, we use a particularly simple QL update mechanism (the one also used for estimates in \cite{VanHasselt:2011}). First explore all actions $N$ times, then propagate to $(s_0,\text{"left"})$. In fact, this is nothing but distributional QL with cyclic exploration and target-matrix trick \cite{dqn} as we explain in Appendix \ref{appendix:theory}. The obtained estimate of $\eta^*(s_0,\text{"left"})$ will be denoted by $\hat \eta(s_0,\text{"left"})$.
\begin{proposition}\label{prop3}
    The sample variance of $\hat \eta(s_0,\text{"left"})$, analogously for "right", after $Nk_1$ steps is $\frac{\sigma_1^2 }{N-1}\chi_{N-1}^2$-distributed. The expectation is $\sigma_1^2$, the variance is $\frac{2\sigma_1^4}{N-1}$.
\end{proposition}
A proof is given in Appendix \ref{appendix:theory}. It is surprising that the sample variance distribution can be identified explicitly as chi-squared, unlike the sample expectations (maxima of independent Gaussians). This is a consequence of the well-known fact in statistics that sample variances of Gaussians are independent of sample expectations. Thus, if return estimates are Gaussian the max-operation of QL (with respect to sample averages) is only delicate for sample expectations, not for sample variances. We emphasize that in contrast to QL the agent in distributional QL does have access to the $\frac{\sigma^2 \chi_{N}^2}{N}$-distributed sample variance by computing sums of the atoms! Hence, the agent can make use of an unbiased estimate for the aleatoric uncertainty $\sigma^2$ which is conflicted by epistemic uncertainty that decreases with $N$. 

\textbf{Implications from of our theoeretical considerations:} We propose the following locally adaptive overestimation mitigation method. (i) Use distributional QL. (ii) At every update compute the sample variance of the current return estimate $\eta_t(s,a)$. (iii) If the sample variance is large replace the QL update by another update (e.g. DQL or an ensemble update or a mixture). Since "large" has no absolute meaning in RL we will compare variances among all actions in $s$ and reduce according to the relative sample variance. 

\textbf{Exploration and overestimation control with ensembles and double $Q$-learning:} 
Known algorithms that directly try to better estimate the $Q$-values require to set up a particular algorithmic architecture. Most use an ensemble of $Q$-copies which are then combined by taking minima \cite{Maxmin}, averages \cite{EnsembleQ}, or averages with random choices \cite{redq}. Ensemble methods are promising in theory (assuming independent ensembles) but more problematic for deep RL as storage problems force ensembles to be parametrized by the same neural network. The optimal number of copies (a hyperparameter) varies for different environments, some choices work well, others fail. 

Our approach is different. We suggest to use QL when it works well (small sample variance) and another algorithm where QL fails (large sample variance). We could combine QL with ensemble methods but for this article decided to combine QL with DQL a bit in the spirit of weighted DQL \cite{weightedDoubleQ} but with completely different weights - we compare to weighted DQL in Appendix \ref{app:ablation}. 
For DQL \cite{hasselt2010} two copies $Q^A$ and $Q^B$ are stored. The update mechanism is similar to QL where the matrix to be updated is chosen randomly in every step. The main difference is the target used, matrices $Q^A$ and $Q^B$ are flipped:
\begin{align*}
     &\quad Q^{A/B}(s,a)\\
   & \leftarrow(1-\alpha)Q^{A/B}(s,a)+\alpha\big(r+\gamma  Q^{B/A}(s',z^*)\big),
\end{align*}
with $z^*=\max_{a'} Q^{A/B}(s',a')$. Here and in the following we use  $Q^{A/B}$ to allow either the choice of $Q^A$ or $Q^B$. Double $Q$-learning reduces the overestimation strongly but sometimes leads to severe underestimation.


The use of sample variance of estimated return distributions is not new to RL. For bandits the simplest example is the UCB exploration bonus for unknown variances that uses this for exploration. In the context of exploration in RL uncertainty dependent exploration has been applied using distributional RL (see for instance \cite{DRL_exploration}, \cite{DRL_exploration2}). The use of sample variances in distributional RL to locally mitigate the overestimation is new to the best of our knowledge. In order to not mix up effects we stick to the standard exploration choices.





\section{ADDQ: Tabular setting}
Based on the theoretical insight we will now introduce a concrete method. We use DQL-updates as an alternative to QL-updates when the agent expects QL-updates to be harmful (large estimated uncertainty by means of sample variance). The weighted DQL-approach we suggest is readily implemented into existing DRL implementations.


\subsection{Weighted DQL: From SARSA-trick to ADDQ}\label{sec:SARSA}

To derive the algorithm let us recall the SARSA convergence proof of \cite{singh2000} that was also used to prove convergence for DQL \cite{hasselt2010} and variants such as clipped $Q$-learning \cite{fujimoto2018addressing} or Maxmin \cite{Maxmin}. The idea is to add a clever $0$, adding and subtracting what is missing to the QL update, thus, writing the algorithm as QL with a bias. If the bias can be proved to disappear, a comparison to QL implies convergence to $Q^*$:
%First recall the tabular DQL update:
% \begin{align*}
% Q^A_{t+1} (s,a) 
%  & = (1-\alpha) Q^A_t(s,a) + \alpha \big( r + \gamma Q^B_t(s^\prime,a^\ast))\\
%   Q^B_{t+1} (s,a)
 % & = (1-\alpha) Q^B_t(s,a) + \alpha \big( r + \gamma Q^A_t(s^\prime,b^\ast))
 %\end{align*}
%where $a^*=\operatorname{argmax} Q_t^A(s,a)$ and $b^*=\operatorname{argmax} Q_t^B(s,a)$.  In contrast to QL two copies of $Q$ are maintained, the update uses the other matrix to determine the maximal entry. Since this entry is typically not identical to the maximal entry of the original matrix the $Q$-values are shifted down compared to QL. The SARSA trick works as follows. Add a clever zero that makes DQL identical to QL plus a bias term:
 \begin{align*}
 Q^{A/B} (s,a) 
  & \leftarrow \overbrace{(1-\alpha) Q^{A/B}(s,a) + \alpha \big( r + \gamma Q^{A/B}(s^\prime,z^\ast)}^{\text{$Q$-learning update}}\\
  &\quad+  \overbrace{\alpha  ( \gamma Q^{B/A}(s^\prime,z^\ast) - \gamma Q^{A/B}(s^\prime,z^\ast) )}^{=:b^{A/B}(s,a)},
 \end{align*}
 with $z^*=\operatorname{argmax} Q^{A/B}(s,a)$. Taking this point of view there are plenty of possibilities to modify DQL to achieve more or less over-/underestimation. As an example, choosing the negative bias-terms 
$b^{A/B}_{\text{clip}}=\alpha \min\{\gamma Q^{B/A}(s^\prime,z^\ast)- \gamma Q^{A/B}(s^\prime,z^\ast),0\}$ yields so-called clipped $Q$-learning introduced as part of TD3 in \cite{fujimoto2018addressing}. As motivated in Section \ref{sec:overestimation} we suggest a locally adaptive overestimation control. The main insight is as follows. One can locally interpolate between QL and DQL by multiplying bias terms with local adaptive weights:
\begin{align*}   
\bar{b}^{A/B}(s,a):=\underbrace{\beta^{A/B}(s,a)}_{\text{new}}\,b^{A/B}(s,a).
\end{align*}
Replacing bias terms $b$ by $\bar b$ generalizes the update. The aggressive choice $\beta=1$ results in QL updates (overestimation), $\beta=0$ results in DQL updates (tendency of underestimation). Choices of $\beta$ suggested in the present article are motivated by the propositions of Sections \ref{sec:overestimationproblem} and \ref{sec:overestimation}. If a lot of uncertainty is present, the algorithm uses large $\beta$, otherwise small $\beta$. The aggressive choices are not necessarily the best, in our experiments below softer choices were more effective. Since overestimation is a priori not problematic for the learning process (if all estimates are equally overestimated the best action does not change, only skewed overestimation slows down the learning) we suggest a choice of $\beta$ that takes into account the local structure, normalizing variances over possible actions.
 
\subsection{New algorithm: locally adaptive distributional double $Q$-learning}\label{sec:tabular}
Following the ideas above we introduce ADDQ, integrating locally adaptive overestimation mitigation in DQL. The pseudo-code given in Algorithm \ref{alg:ADDQ} extends distributional RL pseudo-code to include the double algorithm and adaptive weights. For the tabular target update we follow the measure-mixture approach of \cite{rowland18}, Section 8. Changes to the code for other target updates (for instance gradient steps to minimize KL loss) are straight forward.
\begin{algorithm}[h]
\caption{ADDQ update step}
\label{alg:ADDQ}
\begin{algorithmic}
\REQUIRE Proxies $\eta^A$, $\eta^B$ for $\eta^*$, pair $(s,a)$ to be updated
\STATE Determine step-size $\alpha$
\STATE Sample reward/next state $(r, s')$
\STATE Randomly choose Update(A) or Update(B)
\IF{Update(A)}
\STATE {\color{gray} \# Compute target with locally adapted weight}
\STATE $a^* \gets \operatorname{argmax}_a \mathbb{E}_{Z \sim \eta^A(s',a)} [Z]$
\STATE Determine weight $\beta\in[0,1]$ based on $\eta^A_{old}, \eta^B_{old}$
\STATE $\nu\gets (1-\beta)\eta^B(s',a^*)+\beta\eta^A(s',a^*)$
\STATE $\hat{\eta}_*^A \gets b_{r, \gamma} \#\nu$\\
\STATE {\color{gray} \# Project target back onto support}
\STATE $\hat{\eta}^A \gets \Pi_\mathcal F (\hat{\eta}_*^A)$
\STATE {\color{gray} \# Move $\eta(s,a)$ towards the target, for tabular RL e.g.}
\STATE $\eta^A(s,a) \leftarrow (1 - \alpha)\eta^A(s,a) + \alpha\hat{\eta}^A$
%\STATE $\eta_{new}^A(\bar s,\bar a)=\eta_{old}^A(\bar %s,\bar a)$ for all 
%$(\bar s,\bar a)\neq(s,a)$\\
 \ENDIF
\IF{Update(B)}
\STATE Proceed analogously with $A$ and $B$ exchanged\\
%\OUTPUT $\eta^B_{t+1}$
\ENDIF
\end{algorithmic}
\end{algorithm}
The algorithm is seen to be a combination of distributional QL and distributional DQL. Keeping $\beta$ constant to $1$ is distributional QL, constant $0$ distributional DQL. The key is to chose $\beta$ dependent on the uncertainty that drives the skewed QL overestimation for different actions. 
%\begin{algorithm}[h]
%\caption{ADDQ: tabular categorial setting}
%\label{alg:CDRL}
%\begin{algorithmic}
%\REQUIRE $\eta_t^{A}$, $\eta_t^{B}$ and pair $(s_t,a_t)$ to be updated
%\STATE Determine step-size $\alpha$
%\STATE Sample reward/next state $(r, s')$
%\STATE Randomly choose Update(A) or Update(B)
%\IF{Update(A)}
%\STATE $a^* \gets \operatorname{argmax}_a \mathbb{E}_{Z \sim \eta_t^A(s',a)} [Z]$
%\STATE Determine weight $\beta\in[0,1]$ based on $\eta^A_t, \eta^B_t$% for example based on $w_1(\eta_t^A(S_{t+1},a^*),\eta_t^B(S_{t+1},a^*))$
%\STATE {\color{gray} \# Compute target with locally adapted weight}
%\STATE $\nu\gets \beta\eta_t^B(s',a^*)+(1-\beta)\eta_t^A(s',a^*)$
%\STATE $\hat{\eta}_* \gets b_{r, \gamma} \#\nu$\\
%\STATE {\color{gray} \# Project target back onto support}
%\STATE $\hat{\eta} \gets \Pi_C (\hat{\eta}_*)$
%\STATE {\color{gray} \# stochastic approximation update towars target}
%\STATE $\eta_{t+1}^A(s_t,a_t) = (1 - \alpha)\eta^A_t(s_t,a_t) + %\alpha\hat{\eta}$
%\STATE $\eta_{t+1}^A(s,a)=\eta_{t}^A(s,a)$ for all $(s,a)\neq(s_t,a_t)$\\
%\OUTPUT $\eta^A_{t+1}$
 %  \ENDIF
%\IF{Update(B)}
%\STATE Proceed analogously with $A$ and $B$ exchanged\\
%\OUTPUT $\eta^B_{t+1}$
%\ENDIF
%\end{algorithmic}
%\end{algorithm}
Extending arguments from the literature, notably the convergence proof of \cite{rowland18} for categorical $Q$-learning with stochastic approximation target upate and the SARSA trick of \cite{singh2000}, we prove convergence of Algorithm \ref{alg:ADDQ} for categorical measure parametrizations:
\begin{theorem}\label{thm:double categorical control1}
    Given some initial return distribution functions $\eta_0^{A},\eta^{B}_0$ supported within $[\theta_1,\theta_m]$. If
    %%% TODO: add "arbitrary" (see rebuttal promise)
\begin{itemize}
	\item rewards are bounded in $[R_{min},R_{max}]$ and it holds $[\frac{R_{min}}{1-\gamma},\frac{R_{max}}{1-\gamma}]\subseteq[\theta_1,\theta_m],$
    \item step-sizes fulfill the Robbins-Monro conditions and $\eta^A$ or $\eta^B$ are updated randomly,
    \item the sequences $(\beta^A_t)_{t\in\N},(\beta^B_t)_{t\in\N}$ only depend on the past and fulfill $\lim_{t\to\infty} |\beta^A_t-\beta^B_t|=0$ almost surely,
\end{itemize}
then the induced $Q$-values converge almost surely towards $Q^*$. 
If additionally the MDP has a unique optimal policy $\pi^*,$ then $(\eta_t^{A}),(\eta^{B}_t)$ converge almost surely in $\bar\ell_2$ to some  $\eta^*_C\in\cF_{C,m}$ and the greedy policy with respect to $\eta^{*}_C$ is $\pi^*$.
\end{theorem}

%PLATZHALTER FUER FEHLENDEPLOTS 
\begin{figure*}[h]\label{fig:tabular2}
\begin{center}
\label{g6}
  \includegraphics[height=5cm]{include/Ende.png}
\end{center}
\vspace{-5mm}
\caption{Grid world. First column: Biases summed over all state-action pairs. Other columns highlight the relation of sample variance (bottom) and the effect on $Q$-value (top) estimation. For more experiments and comparisons to Maxmin, EBQL, REDQ see Appendix \ref{app:gred}.}
\vspace{-4mm}
\end{figure*}
%PLATZHALTER FUER PLOTS 



According to Theorem \ref{thm:double categorical control1} symmetric sequences $\beta^A=\beta^B$ that can depend on past distributions $\eta^A, \eta^B$ yield convergence. A particularly simple choice compares the deviations in $\eta^{A/B}$ locally as a local measure for uncertainty.





\textbf{An exemplary choice of adaptive $\beta$:} As motivated earlier QL suffers from overestimation in the presence of uncertainty (function approximation, aleatoric randomness, epistemic randomness) which is reflected in the sample variance (or other measures of the spread of the distribution) of the discrete (random) distributions $\eta^{A/B}_t$. As uniform overestimation is not particularly troubling (e.g. adding a constant to all $Q$-values does not harm at all) the learning process is slowed down by differences in sample variances for allowed actions. There are many other possibilities to choose $\beta$ but we decided to fix this example for all our experiments. For a finite atomic measures $\nu=\sum_{i=1}^n p_i \delta_{a_i}$ define the sample mean $M(\nu)=\frac{1}{n}\sum_{i=1}^n a_ip_i$ and the sample variance $S^2(\nu)=\frac{1}{n-1}\sum_{i=1}^n p_i (a_i- M(\nu))^2$. Now define
\begin{align*}
    S^2_{s,a}&:=\frac{1}{2}\big(S^2(\eta_t^A(s,a)+S^2(\eta_t^B(s,a)\big),\\
    S^2_s&:=\frac{1}{|\mathcal A|}\sum_{a\in \mathcal A}S^2_{s,a},\quad\text{and}\quad
    S^2_{rel}(s,a) := \frac{S^2_{s,a}}{S^2_{s}}.
\end{align*}
According to our computations in the two-sided bandit model, evenly distributed relative sample variances (i.e. all values around $1$) correspond to balanced overestimation effects. Otherwise, overestimation is unbalanced. We define
%
%\begin{enumerate}
 %   \item  Let $V(s,a):=\frac{1}{2}(S^2(\eta_t^A(s,a))+S^2(\eta_t^B(s,a)))$ and 
    %$V_s:=\frac{1}{|\mathcal A|}\sum_{a\in \mathcal A} V(s,a)$.
  %  \item Define relative deviations $V_{rel}(s,a) := \frac{S^2(\eta_t^{A/B}(s,a))}{V_{s}}$. Values close to $1$ (resp. $0$) correspond to relatively large (resp. small) uncertainty among possible actions.    
locally adaptive weights for the next update at $(s,a)$:
\begin{align}\label{adap_rule}
         \beta:= 
        \begin{cases}
        0.75&: S^2_{rel}(s,a) < 0.75 \\
        0.5&: S^2_{rel}(s,a) \in [0.75,1.25] \\
        0.25&: S^2_{rel}(s,a) > 1.25
        \end{cases}.
    \end{align}
    For algorithmic simplicity one could also use squared deviations to the median instead of the mean as the median can be red off directly for measures in the $\mathcal F_{C,m}$ and $\mathcal F_{Q,m}$. The choice combines Q and DQ rather softly, more aggressive choices reduce the length of the interval and increase (resp. decrease) towards $\beta=1$ and $\beta=0$. In practice, (tabular, all Atari, all MuJoCo) the choice worked well without any tuning. The thresholds in the definition of $\beta$ can be seen as hyperparameters to the algorithm. We leave the development of adaptive thresholds for future research.


\subsection{A grid world example}
To show the advantages of ADDQ we present a grid world that highlights, in a more complicated and more realistic way, the main effects of the bandit MDP. For more details, see Appendix \ref{app:gred}. In particular, in Appendix \ref{sec:comparison} we compare ADDQ to other algorithms for overestimation reduction.
\begin{floatingfigure}[r]{2cm}\label{gridword}
\vspace{-3mm}
\begin{center}
\begin{tikzpicture}
\fill[gray!15] (-1.5,-0.5) rectangle (-1,-1);
\fill[gray!15] (-2,1) rectangle (-1.5,0.5);
\fill[gray!15] (-0.5,1) rectangle (0,0.5);
\fill[gray!50] (-1,-1) rectangle (0,0);
\draw[step=0.5cm,color=gray] (-2,-1) grid (0,1);
\node at (-1.75,+0.75) {F};
\node at (-1.25,+0.75) {1};
\node at (-0.75,+0.75) {2};
\node at (-0.25,+0.75) {S};
\node at (-1.75,+0.25) {4};
\node at (-1.25,+0.25) {5};
\node at (-0.75,+0.25) {6};
\node at (-0.25,+0.25) {7};
\node at (-1.75,-0.25) {8};
\node at (-1.25,-0.25) {9};
\node at (-0.75,-0.25) {10};
\node at (-0.25,-0.25) {11};
\node at (-1.75,-0.75) {12};
\node at (-1.25,-0.75) {G};
\node at (-0.75,-0.75) {14};
\node at (-0.25,-0.75) {15};
\end{tikzpicture}
\end{center}
\vspace{-3mm}
\caption{Start in S, goal in G, fake goal in F, high stochasticity and low reward area in gray}
\end{floatingfigure}
Deep RL experiments presented below use essentially  deterministic environments with uncertainty mainly from function approximations. We thus provide a grid world with complicated stochasticity that, depending on parameters, is hard for QL and DQL. The high stochasticity region with low rewards confuses the QL agent. The agent strongly overestimates suboptimal $Q$-values that lead to the gray area, in the gray area $Q$-values are overestimated for actions that stay in the gray area. Hence, all $\varepsilon$-greedy exploration mechanisms spend much time in the gray area. On the other hand, DQL is motivated by estimators of iid random variables and underestimates strongly if action-values for different actions are unevenly distributed. Thus, the DQL agent gets  confused by the local deviations caused by the fake goal and the stochastic region. What results in overestimation for QL, results in underestimation for DQL. In contrast, ADDQ locally combines update-rules of QL and DQL to reduce the estimation biases a lot. In Appendix \ref{app:gred} we show experimentally that the choice of thresholds in $\beta$ is relatively harmless, more or less  aggressive updates still perform well. In contrast, constant $\beta$ and the non-distributional choice from \cite{weightedDoubleQ} with $c=10$ have larger biases.


Plots in the figure above show a few key insights, more in Appendix \ref{app:gred}. First, the estimation bias of ADDQ is much smaller than those of QL and DQL. Most importantly in the complicated states 1, 4, 6, 7. The reason ADDQ better estimates the $Q$-values is the adaptive choice to prefer QL or DQL, depending on (relatively) large variances.



%A number of simple examples has been used to study the effects of algorithms addressing the shortcomings QL. Those examples do not fully represent the problems of QL but  produce convincing plots. Unfortunately, small changes in the parameters strongly change the message of the analysis.  Most prominently, a simple chain MDP used to compare variants of $Q$-learning, see for instance \cite{Sutton1998}, Section 6.7, or \cite{Maxmin}. A more challenging chain MDP was introduced and studied in \cite{EnsembleQ}. According to means positive or negative the example can be used to either support QL or DQL. Another example is simple grid world as used in \cite{hasselt2010}. For the reward variances chosen in \cite{hasselt2010} DQL clearly outperforms QL, while for smaller variances the opposite is the case. In what follows we introduce two types of example MDPs that we believe better capture the problems of QL: chain MDPs and grid world MDPs. They are set up so that situations can be captured in which over- and/or underestimation can be useful or harmful. Most importantly, we provide examples in which QL and DQL both fail. 

%\textbf{Chain MDP:} Figure \ref{figure1} shows our chain MDP example. The MDP is always started in state $A$ and terminates in grey boxes. In state $A$ there are three possible actions left/right/down with zero direct reward. From states $B$, $C$, $D$ there are a number of actions with deterministic next state and Gaussian rewards. The optimal action in $A$ maximizes the true $Q$-values $\mu_B$, $0$, $\mu_C+\gamma \mu_D$.
%\begin{figure}[h]\label{figure1}
%\begin{center}
%\begin{tikzpicture}[scale=0.8]
 %sutton
 %    \node[circle, draw] (A) at (6,0) {A};
  %  \node[text width=0.5cm] at (7,0.5) {0};
   % \node[text width=0.5cm] at (6.5,-0.6) {0};
 %   \node[text width=0.5cm] at (5.3,0.5) {0};
  %  \node[text width=1.7cm] at (3,-1) {$\mathcal N(\mu_B,\sigma^2_B)$};
   % \node[text width=1.5cm] at (8.8,-1) {$\mathcal N(\mu_C,\sigma^2_C)$};
    %\node[text width=1.6cm] at (11,-1) {$\mathcal N(\mu_D,\sigma^2_D)$};
    %\node[circle, draw] (B) at (4,0) {B};
    %\draw[->] (A) -- (B);
    %\node[circle, draw] (C) at (8,0) {C};
    %\draw[->] (A) -- (C);
    %\filldraw[fill=gray, draw=black] (5.75,-1.5) rectangle +(0.5,0.5);
    %\draw[->] (A) -- (6,-1);
    %\filldraw[fill=gray, draw=black] (1.75,-0.25) rectangle +(0.5,0.5);
    %\draw[->] (B) to[out=-105, in=-75]  (2.25,0);
    %\draw[->] (B) to[out=105, in=75]  (2.25,0);
    %\draw[->] (B) to[out=135, in=45]  (2.25,0);
    %\draw[->] (B) to[out=-135, in=-45]  (2.25,0);
    %\fill[black] (3.2,0.25) circle (0.05);
    %\fill[black] (3.2,0) circle (0.05);
    %\fill[black] (3.2,-0.25) circle (0.05);
    %\node[circle, draw] (D) at (10,0) {D};
    %\draw[->] (C) to[out=-75, in=-105]  (9.6125,0);
    %\draw[->] (C) to[out=75, in=105]  (9.6125,0);
    %\draw[->] (C) to[out=45, in=135]  (9.6125,0);
    %\draw[->] (C) to[out=-45, in=-135]  (9.6125,0);
    %\fill[black] (8.8,0.25) circle (0.05);
    %\fill[black] (8.8,0) circle (0.05);
   % \fill[black] (8.8,-0.25) circle (0.05);
   % \filldraw[fill=gray, draw=black] (11.75,-0.25) rectangle +(0.5,0.5);
    %\draw[->] (D) to[out=-75, in=-105]  (11.75,0);
%    \draw[->] (D) to[out=75, in=105]  (11.75,0);
   %\draw[->] (D) to[out=45, in=135]  (11.75,0);
   % \draw[->] (D) to[out=-45, in=-135]  (11.75,0);
   % \fill[black] (10.8,0.25) circle (0.05);
   % \fill[black] (10.8,0) circle (0.05);
   % \fill[black] (10.8,-0.25) circle (0.05);
%\end{tikzpicture}
%\end{center}
%\caption{A chain MDP example}
%\end{figure}
%Now suppose all three means are small and, for instance, $\mu_B<0<\mu_C+\gamma \mu_D$. If $\sigma_B^2$ is sufficiently large, then the overestimation of learning algorithms leads the agent to believe action left would be optimal. Similarly, the underestimation of learning algorithms leads the agent to believe action down would be optimal. Thus, the example is set up in a way that $Q$-learning and variants that do not take into account local properties of the environment and learning uncertainty are doomed to converges slowly.







\newpage
\subsection{ADDQ adaptation for distributional DQN}\label{sec:DQN}
Our DRL based local adaptive overestimation mitigation can be integrated into existing code for DRL with a few extra lines. In the following we describe the integration into C51 \cite{bellemare2017}. We run experiments on Atari environments from the Arcade Learning Environment \cite{ale} using the Gymnasium API \cite{gymnasium}. Algorithms are few line modifications based on the RL Baselines3 Zoo \cite{rl-zoo3} training framework, without any further tuning of hyperparameters.

The C51 algorithm obtained its name from using a categorical representation of return distributions with $m=51$ atoms. The weights of the parametrization are parametrized via feedforward neural networks following the DQN architecture \cite{dqn}. The state $s$ serves as input and the last layer outputs $m=51$ logits for each action followed by a softmax to return probability weights. 
%\begin{floatingfigure}[r]{7cm}
%\begin{center}
%\vspace{-5mm}
%\mbox{
%\includesvg[width=0.5\textwidth]{c51_one}
%}
%\vspace{-6mm}
%\caption{Ex. Asterix, more in Appendix \ref{sec:AC51}}\label{fig:1}
%\end{center}
%\end{floatingfigure}
We write $\eta_\omega(s,a)=\sum_{i=1}^mp_i(s,a;\omega)\delta_{\theta_i}$, 
where $\omega$ comprises the online networks weights. We add a bar $\bar \eta$ to denote a delayed target network which is kept constant and is overwritten from $\eta$ every e.g. 10000 steps with the parameters from the online network. The corresponding expectation is denoted by $Q_\omega(s,a)=\sum_{i=1}^mp_i(s,a;\omega)\theta_i.$ Given a transition $(s,a,r,s')$  the projected target is
\begin{align*}
\hat\eta
=\Pi_C(b_{r,\gamma}\#\eta_{\bar\omega}(s',a^*))=:\sum_{i=1}^m\hat p_i\delta_{\theta_i},
\end{align*} 
where $a^*=\operatorname{argmax}_{a'}Q_{\bar\omega}(s',a')$. Gradient descent on the weights with respect to some loss (here: cross-entropy) is used to move the distribution $\eta_\omega(s,a)$ towards the target distribution $\hat \eta$. As for the tabular algorithm variants with overestimation reduction intervene in the target definition. We keep track of two independently initialized online networks denoted by $\omega^A, \omega^B$ and a pair of respective target networks $\bar\omega^A, \bar\omega^B$. For each gradient step we simulate a vector of random variables with the same size as the batch size with each element determining which of the two estimators is being updated based on the respective transition with the same position in the batch. Accordingly, we use twice the batch size for these methods, so that on average per gradient step, the same number of transitions is used for each estimator, compared to the single-estimator case. We can now describe how to modify the targets for a given transition $(s,a,r,s')$. With the placeholder $\Gamma$ in 
\begin{align*}
    \bar\eta^{A/B}=\Pi_C(b_{r,\gamma}\# \Gamma^{A/B})
\end{align*}
only the place holder is modified for different algorithms:

\textbf{Double C51:} Set $\Gamma^{A/B} = \eta_{\bar\omega^{B/A}}(s',z^*)$ with  $z^*=\operatorname{argmax}_{a'}Q_{\bar\omega^{A/B}}(s',a')$. 

%\textbf{(ii) Double estimator based on target network}, inspired by \cite{hasselt2015doubledqn}: The batch size is the same as in the single estimator case, there are no extra networks. The choice corresponds to 
 %   \begin{align*}
  %      \Gamma &= \eta_{\bar\omega}(s',z^*),\quad z^*=\operatorname{argmax}_{a'}Q_{\omega}(s',a').
   % \end{align*}
    %Greedy action selection is with respect to the online network, i.e. the target network acts as B, the online network as A.\smallskip

   \textbf{Clipped C51}: Inspired by \cite{fujimoto2018addressing}. Set $\Gamma^{A/B}=\eta_{\bar\omega^X}(s',z^*)$ with    $z^*=\operatorname{argmax}_{a'}Q_{\bar\omega^{A/B}}(s',a')$ and $X=\operatorname{argmin}_{c\in\{A,B\}}Q_{\bar\omega^c}(s',z^*)$.\footnote{TD3 [\cite{fujimoto2018addressing}] introduced clipping in an actor-critic setting, where the action is given by the actor. In our C51 adaptation, we select the greedy action based on the target network.}:  

    \textbf{ADDQ (us):} The ADDQ target uses
    \begin{align*}      
    \Gamma^{A/B}=\beta\eta_{\bar\omega^{A/B}}(s',z^*)+(1-\beta)\eta_{\bar\omega^{B/A}}(s',z^*),
    \end{align*}
    where $z^*=\operatorname{argmax}_{a'}Q_{\bar\omega^{A/B}}(s',a').$ The locally adaptive weights $\beta$ may depend on entire state-action return distributions. For the experiments we used the choice from Equation \eqref{adap_rule} based on the online networks $\eta_{\omega^A}, \eta_{\omega^B}$.

%\textbf{(v) Wasserstein reduced mixture:} Denoting by $w_1$ the Wasserstein distance, define
%    \[\Gamma=\sum_{i=1}^m\frac{1}{2}(p_i(s',z^*;\bar\omega^A)+p_i(s',z^*;\bar\omega^B))\delta_{\theta_i-\beta w_1(\eta_{\bar\omega^A}(s',z*),\eta_{\bar\omega^B}(s',z*))},\]
%    where $z^*=\argmax_{a'}(Q_{\bar\omega^A}(s',a')+Q_{\bar\omega^B}(s',a'))$. The negative bias is smaller %the more aligned both distributions are. Since $w_1(\nu,\nu')\leq|\E_{Z\sim\nu}[Z]-\E_{Z\sim\nu'}[Z]|$, %factors fulfilling $\beta\in(0,\frac{1}{2})$ are reasonable. The advantage in comparison to clipping is that %both estimators participate in the target. The concept of measuring the divergence of estimators was also %employed in \cite{weightedDoubleQ}. Now that we have access to the actual distribtuions, this allows us to %measure the estimator's divergence directly based on the Wasserstein distance. Note that for, both the %categorical and the later used quantile representation, there is a simple closed form for $w_1(\nu,\nu').$ We %highlight that this approach can easily be extended to $N>2$ critics for example by subtracting the average %Wasserstein distance from the mixture distribtuion. \smallskip

Results for two Atari environments and a RLiable comparison \cite{RLiable} are presented in Figure \ref{g2} and more extensively in Appendix \ref{app:experiments}. The experiments show that ADDQ is more stable than QL and DQL (it never fails completely) and achieves higher scores in different metrics.

\subsection{ADDQ adaptation for QRDQN} Modifications and results for quantile DRL \cite{dabney2018} are similar to the categorical setting. The setup is explained in Appendix \ref{appendix:QRDQN}; for a quick view two experiments and a RLiable plot for probability of improvements can be found in Figure \ref{g2}. More experimental results are provided in Appendix \ref{app:experiments}. ADDQ is more stable (it never fails completely) and achieves higher scores in different metrics.

\subsection{ADDQ adaptation for quantile SAC}\label{sec:AC}
Compared to Atari environments it is known that overestimation in MuJoCo environments is more severe.  Using the double estimator in critic estimation is not enough, the clipping trick of TD3 \cite{fujimoto2018addressing} greatly improves performance. Recall that, in the formulation of biased $Q$-learning of Section \ref{sec:SARSA}, clipping uses bias $\min\{ Q^{B/A}-  Q^{A/B},0\}$ instead of $Q^{B/A}-Q^{A/B}$. Thus, using our distributional overestimation identification to combine QL and DQL does not hint towards an algorithm competitive with SAC/TD3 or algorithms with more refined overestimation control such as REDQ \cite {redq} or TQC \cite{tqc}. For completeness of the article we still consider the ADDQ effect with standard single and double estimator. We used the base implementations from Stable-Baselines3 to study SAC with quantile regression. With a single critic estimator (we call this QRSAC), with double estimator, clipped double estimator (QB-SAC in the terminology of \cite{tqc}), and ADDQ estimator (ADQRSAC). Experiments are shown in Figure \ref{g3} and more extensively in Appendix \ref{app:experiments}.

As expected ADDQ improves QRSAC and double QRSAC but is not competitive with clipping. We leave for future work to experiment distributional RL based overestimation identification for REDQ and TQC.


%
%
%In continuous control actor-critic methods combine parameterized estimates of the action-value function (critic) and the policy (actor). It is well known that also actor-critic methods tend to suffer heavily from overestimation bias, destabilizing the training process, and reducing final performance. To name but a few ideas 
%\begin{floatingfigure}[r]{7cm}
%\vspace{-5mm}
%\begin{center}
%\mbox{
%\includesvg[width=0.5\textwidth]{mujoco_one}
%}
%\vspace{-7mm}
%\caption{Ex. Humanoid, more in Appendix %\ref{sec:B3}}\label{fig:8}
%\end{center}
%\end{floatingfigure}
%\cite{haarnoja2017,sac,sac_new} use the minimum over two independently initialized critics in the temporal difference target (called clipping). REDQ \cite{redq} minimizes over a randomly selected set of critics (in-target minimization). Truncated quantile critics (TQC) \cite{tqc} is a state-of-the-art extension of SAC using distributional RL. Similarly to distributional DQN in discrete control, the $Q$-function of the critic is parameterized using a quantile network (compare Appendix \ref{appendix:QRDQN}). TQC combines three ideas: distributional repre-
%sentation of a critic, truncation of largest critics atoms, and ensembling of multiple critics. This method significantly increased performance on the MuJoCo benchmark. Our ADDQ adaptation for distributional SAC has a similar effect but controls the overestimation differently. Since the weights as suggested in \ref{adap_rule} cannot be extended directly we use an alternative approach (compare also Appendix \ref{appendix:QRDQN} for the analogue in quantile DQN).
\begin{figure*}[h]\label{fig:experiments}
\begin{center}
  % \label{g1}
  \includegraphics[height=4cm]{include/C51.png}
  \includegraphics[height=2.5cm]{include/C51prob.png}
  %\caption{Our algorithm ADDQ was compared to C51, double C51, and clipped C51 on 10 Atari environments. Here are two exemplary learning curves (averaged over 10 seeds) and the RLiable probability of improvement plot. All details are given in Appendix \ref{}.}
  \label{g1}
  \includegraphics[height=4cm]{include/QRplots.png}
  \includegraphics[height=2.5cm]{include/QRRL.png}
  \vspace{-3mm}
  \caption{Our method ADDQ (blue) compared to  distributional DQN (orange), with double estimator (green) and clipped double estimator (red). First row with categorical, second row quantile parametrization. Learning curves are given for two environments, 8 more in Appendix \ref{app:experiments}. We used 10 seeds. RLiable probability of improvement plots are for all 10 environments, more metrics in Appendix \ref{app:experiments}.}
  \vspace{2mm}\label{g2}
  \includegraphics[height=4cm]{include/SACplot.png}\hspace{-2mm}
  \includegraphics[height=2.5cm]{include/SACRL.png}
  \vspace{-3mm}
    \caption{For completeness: Learning curves for two MuJoCo environments and RLiable probability of improvement on 5 environments. Adapting locally (blue, us) the double critic estimator (green) and direct estimator (orange) cannot (by far) reach performance of the clipping estimator (red). Nonetheless, ADDQ improves distributional direct and double critic estimators with almost no change to the code. Runs are averaged on 10 seeds, more details can be found in the Appendix \ref{app:experiments}.}\label{g3}
  \end{center}
  \vspace{-6mm}
\end{figure*}


%\textbf{Wasserstein Reduced QR-SAC (us):} We build on TQC's framework combining quantile networks with SAC. TQC uses an environment specific number of copies, we always fix two and vary a multiplicative constant $\beta$ (instead of varying the number of copies). The TQC target distributions are replaced by our Wasserstein reduced mixture targets. Additionally, unlike TQC, which updates every critic with every transition drawn from the replay buffer, we randomly select one of the two critics for each transition as suggested by the tabular convergence theorem. This approach prevents too quick an alignment given the identical targets, ensuring that the differences between the two networks arise from more than just their initializations. Given a transition $(s,a,r,s')$, the target distribution is made up from target atoms
%\begin{align*}
%\hat\theta_j &= r + \gamma \Big(\frac{1}{2}(\theta_j(s',a';\bar\omega^A)+\theta_j(s',a';\bar\omega^B))\\
%&\quad-\beta w_1(\eta_{\bar\omega^A}(s',a'),\eta_{\bar\omega^B}(s',a'))-\alpha E\Big)
%\end{align*}
%for $j=1,\dots,m$, where $a'\sim \pi(\cdot:s',\phi)$ is an action given by the current actor network $\phi$, and the deduction of $E=\log(\pi(a';s',\phi))$ comes from the maximum entropy regularization applied in SAC. The bars denote target networks that are occasionally overwritten from the online networks.
%
%Experimental results on MuJoCo environments are presented in Figure \ref{fig:experiments} and more detailed in Appendix \ref{}. We compared five MuJoCo tasks, plotting learning curves and comparisons using RLiable \cite{RLiable} metrics.
%
%
%
%Experiments on MuJoCo environments, Figure \ref{fig:mujoco} of the appendix, show a comparison of  TQC\footnote{In the original paper TQC uses $N=5$ critics. For an equal comparison, we use 2 critics for every algorithm.}. In comparison with a standard (on-adaptive) clipped (resp. double) target and our Wasserstein reduced mixture target. Double can perform well or badly. Adding bias control by clipping improves QR-SAC almost to the level of TQC, our Wasserstein target algorithm (red), even a bit more. It should be mentioned that we did not overengineer the choice of $\beta$, we used $\beta=0.3$ for all environments in addition to Hopper ($\beta=0.5$). This is reasonable, for Hopper also TQC switches to drop 5 instead of 2 atoms per critic.

%The critic is updated by stochastic gradient descent, based on a temporal difference target for policy evaluation similar to the discrete control setting. 
%In this policy evaluation step instead of using the argmax over $a'$ in the target one samples an action from the actor.\smallskip
%
%\textbf{SAC/TD3:} Given a transition $(s,a,r,s')$ from the replay buffer, in the case of SAC [\cite{sac_new}] the update is
%\begin{align*}
%	\gamma \Big(r + Q(s',a';\bar \omega)-\alpha \log(\pi(a';s',\phi)\Big),\quad \text{where}\quad a'\sim \pi(\cdot;s',\phi).
%\end{align*}
%To be precise, the actor network in SAC outputs parameters of a (squashed) Gaussian distribution, from which we can sample the action and calculate the log probabilities. Here the deduction of the log probabilities comes from the maximum entropy regularisation applied in SAC. As this is not the focus of our work, we won't go in to further details concerning the theoretical aspects of the maximum entropy framework  or the update of the actor (policy improvement step) including reparameterization techniques and refer to [\cite{haarnoja2017,sac,sac_new}] instead. Motivated by double $Q$-learning the authors of TD3 \cite{fujimoto2018addressing} proposed to replace $Q(s',a';\bar \omega)$ by the minimum $\min_{X\in\{A,B\}}Q(s',a';\bar \omega^{X})$ over two critic networks $\omega^A$, $\omega^B$ that are alternated in the update mechanism. In the tabular setting this corresponds to the clipping idea recalled in the introduction.
%%The minimization over two separate Q-functions, called clipping, has been proposed by \cite{fujimoto2018addressing}. The authors found that the double estimator [\cite{hasselt2010}] still doesn't sufficiently reduce the overestimation in actor critic methods and therefore introduced the more pessimistic minimization.\smallskip
%
%\todo{Leif: warum brauchen wir eigentlich zwei kopieen und ziehen nicht einfach nur basierend auf einer kopie was ab? wir interpolieren ja nicht zwischen Q und double Q, wen wir irgendwas abziehen}
%\textbf{Adaptive distributional TD3:} Motivated by the adaptive distributional double $Q$-learning ideas from the previous sections we propose the following variant of TD3. We substract some distribution dependent adaptive bias $\beta$ from the $Q$-estimate. ...... MAX ..... Schreiben, was wir tun.
%\smallskip
%
%
%\textbf{Comparison to TQC:}
%\todo{Leif: Auskommentiert. TQC viel kuerzer. Nur in Worten erklaeren, was gemacht wird. Viel wichtiger: erklaeren, warum das anders und schlehcter als unseres ist!}

%TQC [\cite{tqc}] now combines SAC with multiple distributional quantile-estimators based on QRDQN, as critics. The authors introduce a new method to control the overestimation, namely by truncating the topmost quantiles of the mixture distribution in the distributional temporal difference target. Therefore, to compute the target, with $N$ critics as $\eta_{\omega_n}(s,a)=\frac{1}{m}\sum_{i=1}^m\delta_{\theta_i(s,a;\omega_n)}$ for $n\in\{1,\dots,N\}$, TQC first sorts over all the predicted atoms. That is, given a transition $(s,a,r,s')$ and $a'\sim \pi(\cdot;s',\phi)$ denote
%\[z_j(s',a')=sort(\{\theta_i(s,a;\bar\omega_n)| i=1,\dots,m,\ n=1,\dots,N\})_j\]
%where $z_1(s',a')\leq z_2(s',a')\leq\cdots \leq z_{mN}(s',a').$ Then, by means of the hyperparameter $d$, denoting how many atoms per critic to drop, we truncate the topmost $dN$ atoms of the mixture distribution. Therefore, with $k=m-d$ the loss function for the quantile networks (critics) becomes
%\[\cL(\omega)=\frac{1}%{kNm}\sum_{j=1}^{kN}\sum_{i=1}^m\rho_{\tau_i}^1(r+\gamma z_j-\theta_i(s,a;\omega)).\]
%Note that QRDQN sums over the target quantiles while TQC averages.\smallskip



%\begin{figure}[h]
%	\centering
%	\ \includesvg[width=0.95\textwidth]{qrdqn3}
%	\caption{mujoco envs platzhalter}
%	\label{fig:mujoco}
%\end{figure}
\newpage
\section{Summary, limitations, and future work}
Built on theoretical insight for a bandit MDP we suggest to use sample variances in distributional RL to mitigate  the overestimation of QL. Our approach does not use a novel estimation procedure for $Q$-values but instead combines known estimators and tries to use the better one. Using DRL, the agent in state-action pair $(s,a)$ has access to next-state information that predicts the overestimation of QL-updates and accordingly prevers QL- or an alternative. The approach can be incorporated in different estimation methods (e.g. (random) ensemble methods, truncation methods). For the present article we decided to improve DQL, leading to our algorithm ADDQ. The algorithm has the expected feature that in contrast to DQ/DQL it does not fail completely on some environments. Probability of improvement and normalized scores using the RLiable library show clear improvement of ADDQ to underlying deep DQ/DQL.


While ADDQ improves QL and DQL in different settings there is future work to be done. (i) It is clear that our probabilistic calculations can be extended. It is quite likely that results from Gaussian processes can be used for computations in more general settings. (ii) The exemplary choice of $\beta$ can be improved. It would be interesting to replace the hyperparameter thresholds by some adaptive learnable choice.  (iii) The MuJoCo simulation study was only included for completeness, it would have been very surprising if a local adaptation of simple and double critic estimates could improve the clipped critic estimate (or even better  algorithms such as REDQ or TQC). It is interesting future research to include local overestimation mitigation into REDQ (make the chosen ensemble number depend locally on sample variances) or TQC (make the number of truncated atoms depend locally on sample variances). (iv) Use sample variances to perform updates of target networks after non-constant steps.

\section*{Code} The code used in our experiments can be found on GitHub: \url{https://github.com/BommeHD/ADDQ.git}.

\section*{Acknowledgement} The authors acknowledge support by the state of Baden-Württemberg through bwHPC
and the German Research Foundation (DFG) through grant INST 35/1597-1 FUGG.


\section*{Impact statement} This paper presents work whose goal is to advance the field of machine learning, in particular reinforcement learning. There are many potential societal consequences of our work, none which we feel must be specifically highlighted here.

%Some with even more bias than double $Q$-learning (e.g. clipped $Q$-learning from \cite{fujimoto2018addressing}) others with less additional negative bias. Weighted double $Q$-learning \cite{weightedDoubleQ} used a weighted combination of $Q$- and double $Q$-learning estimators. Ensemble Q-learning and averaged $Q$-learning \cite{averagedQ} take averages of multiple action values, an approach that also reduces variance. A more recent idea is to use more than only two copies and combine those to a single estimator, see \cite{Maxmin}. The number of copies can then be used to find the right amount of over- and underestimation to optimise performance for a given model and to reduce variance. The approach has the drawback that the number of copies must be optimised as an additional hyper-parameter and cannot be adapted during a running training process. Bias-corrected Q-Learning \cite{biascorrected} subtracts a bias-term in order to turn $Q$-learning into an unbiased stochastic approximation algorithm. While the idea is natural the implementation is not simple and the analysis requires strong assumptions on the rewards.\smallskip


 
\bibliography{references}
\bibliographystyle{icml2025}


%%%%%%%%%%%%%%%%%%%%%%%%%%%%%%%%%%%%%%%%%%%%%%%%%%%%%%%%%%%%%%%%%%%%%%%%%%%%%%%
%%%%%%%%%%%%%%%%%%%%%%%%%%%%%%%%%%%%%%%%%%%%%%%%%%%%%%%%%%%%%%%%%%%%%%%%%%%%%%%
% APPENDIX
%%%%%%%%%%%%%%%%%%%%%%%%%%%%%%%%%%%%%%%%%%%%%%%%%%%%%%%%%%%%%%%%%%%%%%%%%%%%%%%
%%%%%%%%%%%%%%%%%%%%%%%%%%%%%%%%%%%%%%%%%%%%%%%%%%%%%%%%%%%%%%%%%%%%%%%%%%%%%%%
\newpage
\appendix
\onecolumn
\section{Theoretical backup from probability theory}\label{appendix:theory}
Theoretical backup for all known overestimation reduction algorithms is rather weak, the update rule of $Q$-learning makes precise computations very difficult. As we discuss below, the stochastic approximation rule requires computations with sums of random variables which is not well compatible with computations of maxima of random variables. Justifications for overestimation reduction algorithms are typically based on qualitative arguments for the estimation bias of expectations $\E[\max_i\{ X_1,...,X_K\}]$ of maxima of independent random variables even though random variables appearing in $Q$-learning maxima are clearly not independent. We will give explicit estimates using probabilistic insights for sums and maxima of independent random variables.

We should make very clear that there are different factors to the overestimation problem that are captured in the algorithmic approach of ADDQ. To provide some rigorous theoretical evidence, we focus in this section on the tabular setting without function approximation error. We refer the reader to the seminal article \cite{Thrun+Schwartz:1993} for some simplified computations on the overestimation effect caused by approximation errors.
\subsection{A lower bound computation for the overestimation bias in an episodic bandit MDP (Proof of Proposition \ref{prop2})}


\begin{figure}[h]\label{figure7}
\begin{center}
\begin{tikzpicture}[scale=0.8]
 %sutton
    \node[circle, draw] (A) at (6,0) {$s_0$};
    \node[text width=0.5cm] at (7,0.5) {0};
    %\node[text width=0.5cm] at (6.5,-0.6) {0};
    \node[text width=0.5cm] at (5.3,0.5) {0};
    \node[text width=2cm] at (3.5,-1) {$k_1$  actions};
    %\node[text width=2cm] at (8.5,-1) {$K_C$ identical};
    \node[circle, draw] (B) at (4,0) {$s_1$};
    \draw[->] (A) -- (B);
    %\node[circle, draw] (C) at (8,0) {C};
    \draw[->] (A) -- (C);
% (A) -- (6,-1);
    \filldraw[fill=gray, draw=black] (1.7,-0.25) rectangle +(0.5,0.5);
    %\draw[->] (B) to[out=-105, in=-75]  (2.25,0);
    %\draw[->] (B) to[out=105, in=75]  (2.25,0);
    \draw[->] (B) to[out=135, in=45]  (2.25,0);
    \draw[->] (B) to[out=-135, in=-45]  (2.25,0);
    \fill[black] (3.2,0.25) circle (0.05);
    \fill[black] (3.2,0) circle (0.05);
    \fill[black] (3.2,-0.25) circle (0.05);
    \node[circle, draw] (C) at (8,0) {$s_2$};
    \draw[->] (C) to[out=45, in=135]  (9.65,0);
    \draw[->] (C) to[out=-45, in=-135]  (9.65,0);

    \filldraw[fill=gray, draw=black] (9.7,-0.25) rectangle +(0.5,0.5);
  \fill[black] (8.7,0.25) circle (0.05);
    \fill[black] (8.7,0) circle (0.05);
    \fill[black] (8.7,-0.25) circle (0.05);
        \node[text width=2cm] at (9,-1) {$k_2$  actions};
    \node[circle, draw] (D) at (6.05,-1.39) {};
    \filldraw[fill=gray, draw=black] (5.8,-1.7) rectangle +(0.5,0.5);
    \draw[->] (A) -- (D);
    \node[text width=0.5cm] at (6.5,-0.8) {0};


        
%    \filldraw[fill=gray, draw=black] (7.6,-0.25) rectangle +(0.5,0.5);
    %\draw[->] (C) to[out=-75, in=-105]  (9.6125,0);
    %\draw[->] (C) to[out=75, in=105]  (9.6125,0);
%    \draw[->] (C) to[out=45, in=135]  (9.6125,0);
 %   \draw[->] (C) to[out=-45, in=-135]  (9.6125,0);
  %  \fill[black] (8.8,0.25) circle (0.05);
   % \fill[black] (8.8,0) circle (0.05);
    %\fill[black] (8.8,-0.25) circle (0.05);
\end{tikzpicture}
\end{center}
\end{figure}

The two-sided bandit MDP from the main text can be analyzed side by side. Without loss of generality we thus study the overestimation of the left-side which is essentially the bandit MDP that appeared quite a bit in the literature on overestimation of $Q$-learning (see  \cite{VanHasselt:2011}, Example 6.7 of \cite{Sutton1998}, or \cite{Maxmin}).



%\begin{figure}[h]\label{figure1}
%\begin{center}
%\begin{tikzpicture}[scale=0.8]
 %sutton
 %   \node[circle, draw] (A) at (6,0) {$s_0$};
    %\node[text width=0.5cm] at (7,0.5) {0};
    %\node[text width=0.5cm] at (6.5,-0.6) {0};
   % \node[text width=0.5cm] at (5.3,0.5) {0};
  %  \node[text width=2cm] at (3.5,-1) {$k$ $\mathcal N(\mu,\sigma^2)$ actions};
    %\node[text width=2cm] at (8.5,-1) {$K_C$ identical};
   % \node[circle, draw] (B) at (4,0) {$s_1$};
   % \draw[->] (A) -- (B);
    %\node[circle, draw] (C) at (8,0) {C};
   % \draw[->] (A) -- (C);
% (A) -- (6,-1);
   % \filldraw[fill=gray, draw=black] (1.7,-0.25) rectangle +(0.5,0.5);
    %\draw[->] (B) to[out=-105, in=-75]  (2.25,0);
    %\draw[->] (B) to[out=105, in=75]  (2.25,0);
    %\draw[->] (B) to[out=135, in=45]  (2.25,0);
    %\draw[->] (B) to[out=-135, in=-45]  (2.25,0);
    %\fill[black] (3.2,0.25) circle (0.05);
    %\fill[black] (3.2,0) circle (0.05);
    %\fill[black] (3.2,-0.25) circle (0.05);
    %\filldraw[fill=gray, draw=black] (7.6,-0.25) rectangle +(0.5,0.5);
    %\draw[->] (C) to[out=-75, in=-105]  (9.6125,0);
    %\draw[->] (C) to[out=75, in=105]  (9.6125,0);
%    \draw[->] (C) to[out=45, in=135]  (9.6125,0);
 %   \draw[->] (C) to[out=-45, in=-135]  (9.6125,0);
  %  \fill[black] (8.8,0.25) circle (0.05);
   % \fill[black] (8.8,0) circle (0.05);
    %\fill[black] (8.8,-0.25) circle (0.05);
%\end{tikzpicture}
%\caption{Simple bandit MDP}
%\end{center}
%
%\end{figure}
%The MDP starts in state $s_0$ and goes left (here we skip the right side of the two-sided bandit MDP). Without loss of generality the agent receives reward $0$, the reward could also be chosen deterministic or random without changing the analysis. In state $s_1$ the agent can chose from a number of $k$ actions that all lead to a terminating state and a random reward. The reward distributions are equal. If we denote the reward expectation by $\mu$ then the optimal $Q$-values in $s_0$ are $Q^*(s_0,\text{"left"})=\gamma \mu$ and $Q^*(s_0,\text{"right"})=0$. Thus, if $\mu>0$  the optimal action in state $s_0$ is "left", otherwise "right". In what follows we analyze the bias  $\E[\hat Q(s_0,\text{"left"})]-Q^*(s_0,\text{"left"})$ for an estimator $\hat Q$ obtained from $Q$-learning. Since the action "right" leads to trivial results, without loss of generality we ignore the action "right" and only explore the left hand side of the MDP (otherwise only the notation of exploration steps gets more involved).  

Before stating our results on the overestimation bias we state a technical lemma that is based on the Sudakov-Fernique inequality, see \cite{Sudakov_1971, Sudakov_1976, Fernique_1975}, which is a comparison inequality for Gaussian processes.

\begin{lemma}\label{lemma:martin}
    Let $k \in \mathbb{N}$ and suppose $(X_j^i)_{j \in \mathbb{N}, i \in \{1, \ldots, k\}}$ are iid $\mathcal N(\mu,\sigma^2)$. Further, let $n_{1}, \ldots, n_{k} \in \mathbb{N}$ such that $n_{1} + \cdots + n_{k} \leq k n$ and set $\gamma := \max_{1 \leq i \ne i' \leq k} \bigl|1/n_{n_i} + 1/n_{i'} - 2/n \bigr|$. Then,
    %
    \begin{align}
        \label{eq:diff:max:estimate}
        \Biggl|
            \E\biggl[
                \max_{1 \leq i \leq k} \frac{1}{n_{i}} \sum_{j=1}^{n_{i}} X^i_j
            \biggr]
            -
            \E\biggl[
                \max_{1 \leq i \leq k} \frac{1}{n}\sum_{j=1}^{n} X^i_j
            \biggr]
        \Biggr|
        \leq 
        \sqrt{2 \gamma \ln k}.
    \end{align}
    %
    In particular, if, additionally, $1/n_{i} + 1/n_{i'} \geq 2/n$ for all $i, i' \in \{1, \ldots, k\}$ with $i \ne i'$ then
    %
    \begin{align}
        \label{eq:comparison:max}
        \E\biggl[
            \max_{1 \leq i \leq k} \frac{1}{n}\sum_{j=1}^{n} X^i_j
        \biggr]
        \leq 
        \E\biggl[
            \max_{1 \leq i \leq k} \frac{1}{n_{i}}\sum_{j=1}^{n_{i}} X^{i}_{j}
        \biggr].
    \end{align}
\end{lemma}
%
\begin{proof}
    Given $n_{1}, \ldots, n_{k} \in \mathbb{N}$ with $n_{1} + \cdots n_{k} \leq k n$, set $Y^{i} := n^{-1} \sum_{j=1}^{n} X_{j}^{i}$ and $Z^{i} := n_{i}^{-1} \sum_{j=1}^{n_{i}} X_{j}^{i}$ for any $i \in \{1, \ldots, k\}$. Clearly, $Y = (Y^{1}, \ldots, Y^{k})$ and $Z = (Z^{1}, \ldots, Z^{k})$ are Gaussian vectors with $\E[Y^{i}] = \E[Z^{i}]$ for all $i \in \{1, \ldots, k\}$ and
    %
    \begin{align*}
        \gamma_{i, i'}^{Y}
        :=
        \E\Bigl[\bigl(Y^{i} - Y^{i'}\bigr)^2 \Bigr]
        =
        2/n
        \qquad \text{and} \qquad
        \gamma_{i, i'}^{Z}
        :=
        \E\Bigl[\bigl(Z^{i} - Z^{i'}\bigr)^2 \Bigr]
        =
        1/n_{i} + 1/n_{i'}
    \end{align*}
    %
    for all $i, i' \in \{1, \ldots, k\} $ with $i \ne i'$. Thus, \eqref{eq:diff:max:estimate} is an immediate consequence of \citealp[Theorem~2.2.5, Eq.~(2.2.11)]{Adler_2007}. If, in addition, $1/n_{i} + 1/n_{i'} \geq 2$ for all $i \ne i'$ then $\gamma_{i, i'}^{Y} \leq \gamma_{i, i'}^{Z}$ and \eqref{eq:comparison:max} follows from \citealp[Theorem~2.2.5, Eq.~(2.2.12)]{Adler_2007}.
\end{proof}

We are now in a position to give a lower bound on the overestimation bias for all exploration policies under the typical $\frac{1}{\# \text{visits}}$-step-size schedule. This extends the overestimation analysis of \cite{VanHasselt:2011} which was based on a simple synchronous exploration mechanism. The step-size schedule is crucial for our probabilistic analysis as it turns the analysis in a computation with sums and maxima of independent random variables. The sum structure is a consequence of the simple fact that $z_t:=\frac{1}{t} \sum_{i=1}^t a_{i-1}$ solves the stochastic approximation recursion $z_{t+1}=(1-\alpha_t)z_t+ \alpha_t a_t, z_0=0,$ with $\alpha_t=\frac{1}{t+1}$. Here is a formal lower bound of the overestimation bias in the episodic bandit MDP. As mentioned above we only consider exploration of the left-side of the bandit MDP.
\begin{theorem}\label{thm:1}
    Suppose $Q_0$ is initialized as the zero matrix and step-sizes are chosen as  $\alpha_t(s,a)=\frac{1}{T_{s,a}(t)}$, with $T_{s,a}(t)$ the number of updates of $Q(s,a)$. If rewards are $\mathcal N(\mu,\sigma^2)$-distributed, then every sufficient exploratory exploration rule leads to overestimation bias at least 
    \begin{align*}
        \E[Q_{Nk}(s_0,\text{"left"})]- Q^*(s_0,\text{"left"})\geq\frac{\gamma}{\sqrt{\pi\log(2) }}\frac{\sigma \sqrt{\log(k)}}{\sqrt{N}}.
    \end{align*}
    By sufficient exploratory we mean that $\frac{1}{n_a(t)}+\frac{1}{n_{a'}(t)}
\geq \frac{2}{N}$ for all actions $a,a'$ with $n_a(t)=T_{s_{,a}}(t)$.
\end{theorem}
In simple words, the expected overestimation in $n$ episodes of $Q$-learning on the left-side of the bandit MDP is at least of the order $\frac{\sigma \sqrt{k\log(k)}}{\sqrt{n}}$. Higher variance and more actions obviously lead to stronger overestimation. 

\begin{proof}
    In contrast to \cite{VanHasselt:2011} we do not assume synchronous  exploration of all actions in each episode. Instead, we give a comparison argument that allows to reduce sufficiently exploratory exploration to simple cyclic exploration. The approach is motivated by the target network (here: target matrix) trick from DQN \cite{dqn}. With target networks the update works as follows:
     \begin{align*}
        Q_{t+1}(s,a)\leftarrow (1-\alpha_t) Q_t(s,a)+\alpha_t \big(r+\gamma \max_{a'}\bar Q_t(s',a')\big),
    \end{align*}
    where $\bar Q$ is kept constant for a fixed number of steps and then updated from the current $Q$-matrix. Although the target network was introduced in DQN to reduce overfitting in function approximation, it also reduces the overestimation of $Q$-values. We use the latter effect locally in this proof to construct a lower bound. The target matrix is kept constant ($0$ matrix) at $s_0$ until the last update. We show that (i) this target matrix trick with cyclic exploration yields a lower bound for the overestimation bias of standard $Q$-learning with sufficient exploratory  exploration strategy and (ii) allows for computations using elements of probability theory. In what follows we denote by $X_i^a$ the reward obtained when playing chosing $a$ for the $i$th time. By assumption all $(X_i^a)$ are iid.
    
    \textbf{Step 1:} In the first step we show that using cyclic exploration (playing one action after the other) for $N$ rounds ($Nk$ steps in total) followed by one update at $(s_0,\text{"left"})$ minimizes the overestimation of $Q$-learning estimators $\hat Q^{cyc,tar}(s_0,\text{"left"})$ that explore each action at most $N$ times.   
    For the claim we compare arbitrary $Q$-learning with the cyclic variant.    
    Recalling the $Q$-learning update (with arbitrary exploration rule) and the step-size schedule shows that 
    \begin{align*}
        Q_{t}(s_1,a)=\frac{1}{T_{s,a}(t)}\sum_{j=1}^{T_{s,a}(t)} X^a_j\sim \mathcal N(\mu,t\sigma^2).
    \end{align*}
    Since $s_0$ is explored once per episode the step-size schedule of regular $Q$-learning yields
    \begin{align*}
        Q_{Nk}(s_0,\text{"left"})= \frac{1}{Nk}\sum_{t=1}^{Nk} \gamma\max_a Q_t(s_1,a)
    \end{align*}
    for the $N$th episode. We next show that $\E[Q_{Nk}(s_0,\text{"left"})]\geq \E[\hat Q^{cyc,tar}(s_0,\text{"left"})]$. Denoting $n_k(t)=T_{s_1,a_k}(t)$ and using Lemma \ref{lemma:martin} we obtain
    \begin{align*}
        \E[Q_N(s_0,\text{"left"})]
         &=\E\Big[\frac{1}{Nk}\sum_{t=1}^{Nk} \gamma\max_a\big\{ Q_t(s_1,a_1),...,Q_t(s_1,a_k)\big\}\Big]\\
        &=\frac{1}{Nk}\sum_{t=1}^{Nk}\gamma \E\Big[\max \Big\{\frac{1}{n_1(t)}\sum_{j=1}^{n_1(t)}X^{a_1}_j,...,\frac{1}{n_k(t)}\sum_{j=1}^{n_k(t)}X^{a_k}_j\Big\}\Big]\\
        &\geq\frac{1}{Nk}\sum_{t=1}^{Nk}\gamma \E\Big[\max \Big\{\frac{1}{N}\sum_{j=1}^{N}X^{a_1}_j,...,\frac{1}{N}\sum_{j=1}^{N}X^{a_k}_j\Big\}\Big]\\
        &=\gamma \E\Big[\max \Big\{\frac{1}{N}\sum_{j=1}^{N}X^{a_1}_j,...,\frac{1}{N}\sum_{j=1}^{N}X^{a_k}_j\Big\}\Big]\\
        &=\E[\hat Q^{cyc,tar}(s_0,\text{"left"})].
    \end{align*}
\textbf{Step 2:} We now analyze the overestimation bias for the $Q$-learing with target matrix trick after $Nk$ episodes. To deduce the claim we use a fact on the expectation of the maximum of independent $\mathcal N(\mu,\sigma^2)$-distributed random variables:
\begin{align*}
    \mu+\frac{1}{\sqrt{\pi\log(2) }}\sigma\sqrt{\log(k)}\leq \E[\max\{X_1,...,X_k\}]\leq \mu+\sqrt{2}\sigma\sqrt{\log(k)}
\end{align*}
The inequality can for instance be found in \cite{vH}. According to the step-size schedule the update is as follows. If $N=nk$, i.e. in cyclic exploration every action is played $n$ times, then $Q^{cyc,tar}_{Nk}(s_1,a)=\frac{1}{N} \sum_{i=1}^N X_i^a$ for independent $\mathcal N(\mu,\sigma^2)$-distributed reward samples. Thus, $Q_{Nk}(s_1,a_1),...,Q_{Nk}(s_1,a_k)$ are iid and $\mathcal N(\mu,\frac{\sigma^2}{N})$-distributed. The final update after $N$ episodes is to set $$\hat Q^{cyc,tar}(s_1,\text{"left"})= \gamma \max \{Q_{Nk}(s_1,a_1),...,Q_{Nk}(s_1,a_k)\}\overset{(d)}{=} \gamma \max \{Z_1,...,Z_k\}$$ 
for $Z_1,...,Z_k$ iid $\mathcal N(\mu,\frac{\sigma^2}{N})$. The claim follows from the lower bound on the expectation of maxima of independent Gaussians.
\end{proof}

%We can prove a similar variant for Bernoulli rewards, i.e. rewards taking values $+1$ (resp. $-1$) with probabilities $p$ (resp. $1-p$).
%\begin{proposition}
 %  In the same setting as Theorem \ref{thm:1} with $Ber(p)$-distributed rewards we have, for the number of actions $k$ large,
 %\begin{align*}
  %      \E[Q_{N}(s_0,\text{"left"})]\gtrapprox Q^*(s_0,\text{"left"})+ \sqrt{\frac{2p(1-p) k\log(k)}{\lceil\frac{N}{k}\rceil k}}.
  %  \end{align*}
  %  after $N$ episodes.
  %  \end{proposition}
%\begin{proof}
 %   The proof strategy is equal but the probability theory involved is different. First, Lemma \ref{lem:stochastic approximation lemma} can be proved similarly for Bernoulli random variables. Thus, the cyclic target matrix variant of $Q$-learning provides a lower bound for the overestimation bias. Unfortunately, we are not aware of a convenient lower bound on $\E[\max\{X_1,...,X_k\}]$ of sums of independent Bernoulli variables (i.e. Binomial(N,p) variables). We thus rely on the asymptotic behavior derived from the extremal value centeral limit theorem: 
%    \begin{align*}
 %       \lim_{k\to\infty}\E[\max\{X_1,...,X_k\}]=Np+\sqrt{2p(1-p)\log(k)}.
 %   \end{align*}
%\end{proof}
The computations give a number of quantitative insights. First, it becomes very clear why explicit computations for $Q$-learning are complicated. The stochastic approximation update is closely related to sums (exactly equal to sums for $\frac{1}{\# \text{visits}}$  step-sizes, while the update target involves a maximum. Unfortunately, sums and maxima of random variables do not get along well. We thus performed computations for Gaussian reward distributions for which sums are Gaussian again. Maxima of Gaussian random variables (processes) are a well-studied field in Mathematics and estimates can be derived. 

Finally, the computations show how variances influence the overestimation problem. Stochastic rewards with high variance are more problematic than rewards with small variance. In the next section we study distribution RL in the same simple problem.

\subsection{On the use of variance in local overestimation control via distributional RL (Proof of Proposition \ref{prop3})}
Our computations for the bandit MDP above suggest uncertainty at next states (aleatoric uncertainty of rewards, epistemic uncertainty through small $N$) lead to larger overestimation. Thus, algorithms mitigating the overestimation problem locally at $(s,a)$ must somehow use uncertainty at next states $s'$. In the following proposition we analyze the distributional $Q$-learning update scheme (compare Section 8 of \cite{rowland18})
\begin{align*}
    \eta(s,a)\leftarrow (1-\alpha) \eta(s,a)+\alpha\big(b_{r,\gamma}\#\eta(s',a^*)\big),
\end{align*}
with $a^*=\operatorname{argmax}_{a'}Q(s',a')$, where $Q(s',a')$ are the (random) sample means of the estimated return distributions $\eta(s',a')$. The bootstrap function is $b_{r,\gamma}(z)= r+\gamma z$ and $f\#\nu(B):=\nu(f^{-1}(B))$ denotes the push-forward of measures. Recall from the main text that distributional QL learns random measures, $\eta(s,a)$ itself is a probability measure that depends on the random samples used in the updates. As a probability measure $\eta(s,a)$ has an expectation and a variance that are both random in terms of the random samples. The situation becomes a bit tricky as we are going to take variances of the variances. To avoid confusion we use an analogy to statistics and speak of (random) sample averages $M$ (resp. sample variances $S^2$) of $\eta(s,a)$ and expectation $\E$ (resp. variance $\mathbb V$) for the integrals against the true randomness induced by the probability space behind all appearing random variables.

Motivated by the previous section we only study cyclic exploration with "target measure", i.e. all actions in $s_1$ are explored $N$ times before bootstrapping the estimates to $s_0$. Using the standard $\frac{1}{\# \text{visits}}$-step-size schedule allows us to carry out fairly explicit computations using tools from probability theory. Most importantly, we can derive the exact distribution for the sample variance of $\eta(s_0,\text{"left"})$, the state-action pair at risk for overestimation. The insight obtained from the computation motivates our ADDQ algorithm.
\begin{proposition}\label{prop1}
    Consider the bandit MDP introduced above with $\mathcal N(\mu,\sigma^2)$ reward distributions and $k$ actions. Denote by $\hat{ \eta}^{cyc,tar}(s_0,\text{"left"})$ the return distribution of distributional $Q$-learning with cyclic exploration, step-size schedule $\alpha_t(s,a)=\frac{1}{T_{s,a}(t)}$, and "target distribution". More precisely, all arms are $N$ times explored before the first and only update at $(s_1,\text{"left"})$. It then follows that
    \begin{align*}
            S^2 (\hat{ \eta}^{cyc,tar}(s_0,\text{"left"}))\sim \frac{\sigma^2}{N-1} \chi_{N-1}^2,
    \end{align*}
    where $\chi_n^2$ denotes the chi-squared distribution with $n$ degrees of freedom.


    
\end{proposition}
It is interesting to note that while there is a simple distributional expression for the sample variance there is no simple expression for the distribution of the sample mean $\max\{\frac{1}{N}\sum_{i=1}^N X_i^{a_1},...,\frac{1}{N}\sum_{i=1}^N X_i^{a_k}\}$ which is the maximum of independent Gaussians.
\begin{proof}
    For the proof we need to spell-out the distributional $Q$-learning update and then use basic results from statistic for sample-mean/sample-variances.

    Let us first understand the distributional update at the last step. Similar to scalar stochastic approximation schemes, the $\frac{1}{n+1}$-step-size schedule also gives the distributional estimator
    \begin{align}
        \hat\eta^{cyc, tar}(s_1,a)=\frac{1}{N}\sum_{i=1}^N \delta_{X_i^a}
    \end{align}    
    for the pre-terminal state $s_1$ after $N$ explorations of action $a$. To see why write $\eta_n:=\frac{1}{n}\sum_{i=1}^n\delta_{X^a_i}$ to get
    \begin{align*}
        \eta_{n+1}=\frac{1}{n+1}\sum_{i=1}^{n+1} \delta_{X_i^a}
        =\Big(1-\frac{1}{n+1}\Big)\frac{1}{n}\sum_{i=1}^{n} \delta_{X_i^a}+\frac{1}{n+1}\delta_{X^a_{n+1}}
        =(1-\alpha_n) \eta_n+\alpha_n \delta_{X^a_{n+1}}
    \end{align*} 
    and recall that there is no max-term in the update for pre-terminal states. Note that the equal weights $\frac{1}{N}$ are a consequence of the step-size schedule. To compute the return distribution at $s_0$ for action "left" we need to identify $a^*$. For that sake we must identify the sample means of $\hat \eta^{cyc,tar}(s_1,a)$:
    \begin{align*}
        \hat Q^{cyc,tar}(s_1,a)=\frac{1}{N}\sum_{i=1}^N X_i^a \sim \mathcal N(\mu,\sigma^2/N)
    \end{align*}
    Now we chose $a^*=\argmax_a \hat Q^{cyc,tar}(s_1,a)$ and set    
    \begin{align*}
        \hat\eta^{cyc, tar}(s_0,\text{"left"})
        =b_{0,\gamma} \# \hat\eta^{cyc,tar}(s_1,a^*)
        =: \gamma X.
    \end{align*}
    While we do not know much about $X$ (it is the empirical distribution of the set of Gaussians $X_1^a,...,X_N^a$ with maximal sum) we know the exact distribution of the sample variance for the following reason. The sample variance of every $\eta^{cyc, tar}(s_1,a)$ is nothing what is called sample variance of the iid observation $X_1^a,...,X_k^a$ in statistics. If $M_a:=\frac{1}{N}\sum_{i=1}^N X_i^a$ is the sample mean and $S_a^2:=\frac{1}{N-1}\sum_{i=1}^N(X_i^a-M^a)^2$ the sample variance then it is well-known that 
    \begin{itemize}
        \item $\E[S^2_a]=\sigma^2$,
        \item $S_a^2$ is $\frac{\sigma^2}{N-1}\chi^2_{N-1}$-distributed,
        \item $M_a$ and $S^2_a$ are independent.
    \end{itemize}
    The third property implies that if $S^2_{a_1},...,S^2_{a_k}$ are independent sample variances and $a^*$ is chosen to maximize the sample means, then also $S^2_{a^*}\sim \frac{\sigma^2}{N-1} \chi^2_{N-1}$ while $M_{a^*}\nsim M_a$. 
\end{proof}
As a consequence, even though we cannot compute the distribution of the sample mean of $\hat\eta^{cyc, tar}(s_0,\text{"left"})$ we can compute the distribution of its sample variance as $\frac{\sigma^2}{N-1}\chi^2_{N-1}$. The expectation is $\E[S_a^2]=\sigma^2$ while the variance is $\mathbb V(S_a^2)=\frac{2\sigma^4}{N-1}$. Note that one might be tempted to believe the sample variance is independent of the number $k$ of actions. This is not the case, as the number of explorations $Nk$ needed for $N$ explorations depends on $k$. Alternatively, one might denote the number of total episodes by $n$ and replace $N$ by $\lfloor\frac{n}{k}\rfloor$.


%
%
%
%\begin{proposition}\label{prop1}
 %   Consider the bandit MDP introduced above with $Ber(p)$ reward distributions. Then, after $Nk$ episodes of distributional $Q$-learning with cyclic exploration and step-sizes $\alpha_t(s,a)=\frac{1}{T_{s,a}(t)+1}$ the estimated optimal return distribution is
  %  \begin{align*}
   %    \hat{ \eta}^{cyc,tar}(s_0,\text{"left"})\overset{(d)}{=} \gamma Z,
   % \end{align*}
    %where $Z$ is Bernoulli distributed with (random) parameter $\hat p\overset{(d)}{=}\frac{1}{N} \max\{Z_1,...,Z_k\}$ for iid $Bin(N,p)$ random variables. In particular, for large $k$, $\hat p$ has approximate expectation $p+\sqrt{\frac{2p(1-p) \log(k)}{N}}$.
%\end{proposition}



%\begin{proof}
 %   Let us first understand the distributional update at the last step. Similar to scalar stochastic approximation schemes, the $\frac{1}{n}$-step-size schedule also gives the distributional estimator
  %  \begin{align}
   %     \hat\eta^{cyc, tar}(s_1,a)=\frac{1}{N}\sum_{i=1}^N \delta_{X_i^a}
    %\end{align}    
    %for the pre-terminal state $s_1$ after $N$ explorations of action $a$. To see why write $\eta_n:=\frac{1}{n}\sum_{i=1}^n\delta_{X^a_i}$ to get
   % \begin{align*}
    %    \eta_{n+1}=\frac{1}{n+1}\sum_{i=1}^{n+1} \delta_{X_i^a}
     %   =\Big(1-\frac{1}{n+1}\Big)\frac{1}{n}\sum_{i=1}^{n} \delta_{X_i^a}+\frac{1}{n+1}\delta_{X^a_{n+1}}
      %  =(1-\alpha_n) \eta_n+\alpha_n \delta_{X^a_{n+1}}.
    %\end{align*} 
    %Now we assume all rewards are $Ber(p)$-distributed. In that situations the return distributions only have atoms at $+1$ and $-1$ and it suffices to compute the (random) weights depending on the rewards observed. Counting outcomes $+1$ and $-1$ yields
    %\begin{align*}
     %   \hat\eta^{cyc, tar}(s_1,a)=\frac{\#\{X^a_i=1\}}{N}\delta_1+\frac{\#\{X^a_i=+1\}}{N}\delta_{-1}.
    %\end{align*}
    %In words, the estimated return distributions $\hat\eta^{cyc, tar}(s_1,a)$ are independent Bernoulli-distributed with (random) weights $\hat p_a\sim\frac{1}{N} Bin(N,p)\in [0,1]$.
    
    %Next, we compute the (random) return distribution at $s_0$ for action "left", the action "right" is irrelevant. According to the target matrix update rule and the choice of step-sizes the first (and only) update at $(s_0,\text{"left"})$ yields  
    %  \begin{align*}
     %   \hat\eta^{cyc, tar}(s_0,"left")
      %  =b_{0,\gamma} \# \hat\eta^{cyc,tar}(s_1,a^*)
       % = \gamma X,
    %\end{align*}
    %where $X$ is Bernoulli distributed with (random) parameter $\hat p=\frac{1}{N}\max\{\hat p_1,...,\hat p_k\}\overset{(d)}{=}\frac{1}{N} \max\{Z_1,...,Z_k\}$ for iid $Bin(N,p)$ random variables. The first equality is the distributional $Q$-learning update at a pre-terminal state  with $\alpha=1$ (first update). For the second equality we argue as follows. First, $a^*$ is chosen according to the maximal expectations of the (random) measures $\hat\eta^{cyc,tar}(s_1,a)$. All these measures are Bernoulli with (random) success probability $p$ distributed according to $\frac{1}{N}Bin(N,p)$. 
    
    %The statement on the the expectation of $\hat p$ follows from extremal value theory from the expectation of the maximum of $k$ independent binomial variables for large $k$:
    %\begin{align*}
     %   \E[\max\{Z_1,...,Z_k\}]\approx Np+\sqrt{2p(1-p)N\log(k)}.
    %\end{align*}
%    Now we continue with the $Q$-learning update to obtain $Q(A,left).$ For that sake we need to compute the maximal expectation. The expectations are $p=\frac{1}{n} Bin(n,p)$, so we need the maximum of independent binomial distributions, which is known to behave like
 %   \begin{align*}
  %      1/n(np+\sqrt{2p(1-p}\sqrt{n log(k)})=p+\sqrt{2p(1-p)}\sqrt{k}\frac{1}{\sqrt{n}}
   % \end{align*}
  %  for large $k$. Now writing down the update gives
   % \begin{align*}
    %    \eta_n(s_0,"left")&=(1-\alpha) \delta_0+ \alpha(0+\gamma \eta(C,a^*)\\
     %   &=(1-\alpha)\delta_0+\alpha \gamma X)
    %\end{align*}
    %with $X\sim Bernoulli(p+\sqrt{2p(1-p)}\sqrt{\log k}\frac{1}{\sqrt{n}}$. In particular, we see the overestimation $Q(A,left)=p+\sqrt{2p(1-p)\log k}\frac{1}{\sqrt{n}}$.    
%\end{proof}
%The choice of Bernoulli-rewards were made in order to follow explicitly the distributional update rule. What can be seen  explicitly is the effect of aleatoric and epistemic randomness. The first summand $p$ in the approximate expectation of $\hat p$ is due to the inherent aleatoric randomness induced by the Bernoulli rewards while the second summand $ \sqrt{\frac{2p(1-p)\log(k)}{N}}$ is due to the epistemic randomness of the learning process. As can be expected in this example we can see clearly that aleatoric and epistemic uncertainty is not independent, high aleatoric uncertainty also implies high epistemic uncertainty as the estimation of high-variance random variables requires more effort.


%Since Bernoulli random variables with parameter $p$ have variance $p(1-p)$ we obtain the following corollary:
%\begin{corollary}
 %   In the setting from Proposition \ref{prop1} the (random) variance of the return distribution $\hat\eta^{cyc,tar}(s_0,\text{"left"})$ is $\gamma^2\hat p(1-\hat p)$ with $\hat p\overset{(d)}{=}\frac{1}{N} \max\{Z_1,...,Z_k\}$ for iid $Bin(N,p)$ random variables.
%\end{corollary}
%At the beginning of the learning procedure the variance of $\hat \eta^{cyc,tar}(s_0,\text{"left"})$ is clearly effected by the epistemic contribution $\sqrt{\frac{2p(1-p)\log(k)}{N}}$ which diminishes for increasing $N$.
We now come back to the motivation of our ADDQ algorithm, in particular the choice of $\beta$. Let us consider the two-sided bandit MDP with $\mathcal N(\mu_1,\sigma_1^2)$ (resp. $\mathcal N(\mu_2,\sigma_2^2)$) distributed rewards. Suppose $\mu_1>\mu_2$ are small but $\sigma_1^2\ll \sigma_2^2$  and/or $k_1\ll k_2$. The MDP is delicate as the following can happen using QL (overestimation) and DQL (global overestimation reduction). In QL the agent will believe for a long time that "right" is the optimal action in $s_0$. In DQL the agent will believe for a long time that "down" is the optimal action as both non-trivial $Q$-values can be underestimated to be negative. It is thus more reasonable to mitigate overestimation locally instead of globally. This is what ADDQ achieves through the choice of $\beta$.

To use our explicit computations above, the agent explores both sides with cyclic exploration and target-matrix update. Now assume both sides have been explored with a total number of $n$ episodes each. The agent knows the estimates $\hat \eta(s_0,\text{"left"})$ (resp. $\hat \eta(s_0,\text{"right"})$) and the corresponding $Q$-values (sample means of return distributions) $\hat Q(s_0,\text{"left"})$ (resp. $\hat Q(s_0,\text{"right"})$). From the overestimation study of the previous section, the agent unwittingly overestimated the true $Q$-values in the order of $\frac{\sigma_1 \sqrt{k_1\log(k_1)}}{\sqrt{n}}$ (resp. $\frac{\sigma_2 \sqrt{k_2\log(k_2)}}{\sqrt{n}})$ and will believe "right" is the correct action. Now the smart agent also knows he/she should take into account the sample variances that is known to the agent and we know are of the order $\sigma_1^2$ (resp. $\sigma_2^2$) with concentration depending on $k_1$ (resp. $k_2$). The local overestimation control of $\beta$ (based on relative sample variances) from \eqref{adap_rule} thus compares $\sigma_1^2$ to $\sigma_2^2$ and suggests to mitigate overestimation of $\hat Q(s_0,\text{"right"})$. In our algorithm we use double $Q$-learning to mitigate, the same idea can of course be integrated into other algorithms (such as changing number of truncation atoms in truncated quantile critics algorithms \cite{tqc}).

%For simplicity we again use the target network trick, other exploration and update schedules lead to even stronger overestimation (see last section). According to Proposition \ref{prop1} distributions $\hat \eta(s_0,\text{"left"})$ and $\hat \eta(s_0,\text{"right"})$ are Bernoulli with (random) parameters satisfying $\E[\hat p_1]\approx p_1+\sqrt{\frac{2p_1(1-p_1) \log(k_1)}{N}}$ and $\E[\hat p_2]\approx p_2+\sqrt{\frac{2p_2(1-p_2) \log(k_2)}{N}}$. What can be seen very clearly in this example is the following. By considering the (random) variance (here equivalent to (random) success probability $\hat p$) of return distributions in $(s,a)$ the agent obtains inside on the aleatoric and epistemic uncertainty in $s'$ which again lead to overestimation of $Q$-values. Since both imply overestimation ADDQ reduces overestimation where the relative variance is large. 

\newpage

\section{Experimental confirmation of theoretical results for two-sided bandit MDP}

We provide an experimental analysis of the two-sided bandit MDP for which we proved theoretical results on the overestimation. For reimplementation purposes we collect here all required information on the environment and the training for all plots.

Environment:
\begin{itemize}
    \item $\gamma=0.9$
    \item $\mu_1=-0.1$, $\mu_2=0.1$, $k_2=5$, $\sigma_2=1$; the correct decision is thus moving to the right in the Start State.
\end{itemize}
Distributional properties for ADDQ:
\begin{itemize}
    \item categorical parametrization, 51 atoms equally spaced on $\lbrack-3,3\rbrack$
    \item initialization as $\delta_0$
\end{itemize}
Algorithmic choices:
\begin{itemize}
    \item $\beta$ is chosen according to (\ref{adap_rule})
    \item step-size schedules $\alpha_t(s,a)=\frac{1}{T_{s,a}(t)}$, with $T_{s,a}(t)$ the number of visits in $(s,a)$ up to time $t$, i.e. $\frac{1}{n}$ state-action wise counted 
    \item exploration: either $\epsilon$-greedy with $\epsilon$ linearly decreasing from $1$ to $0.1$ in $10000$ steps, then constant (E) or uniform random (U)
\end{itemize}
Policy evaluation: evaluation of 3 steps with a frequency of every 500 steps using current greedy policy, correct action rates refer to if the exploration was greedy

In order to demonstrate the proven proportionality of the overestimation of $Q$-values by QL in the number of arms and the variance to the left, we keep one of the values fixed and compare different values in the other one.
\begin{itemize}
    \item First experiment: $\sigma_1=5$ fixed, iterating over $k_1=5,10,15,20$, denoted as $K$ in the legend
    \item Second experiment: $k_1=10$ fixed, iterating over $\sigma_1=2,4,6,8$, denoted as $S$ in the legend
\end{itemize}

The plots also demonstrate that

\begin{itemize}
    \item ADDQ is much better in terms of bias, leveraging local information given by the (relative) variances
    \item the variances and relative variances at state 0 capture the real variances given by the MDP
    \item although our theorems assumed a sufficiently exploratory policy, the results seem to generalize to the much more commonly used $\epsilon$-greedy setting
\end{itemize}
\newpage

\begin{figure}[H]
    \centering
    \includegraphics[width=0.95\textwidth]{include/SuttonExtended_K_E_Big.png}
    \includegraphics[width=0.95\textwidth]{include/SuttonExtended_K_U_Big.png}
    \caption{Comparing ADDQ and QL on two-sided bandit MDP with different number of arms on the left side and different exploration settings.}
    \label{fig:11}
\end{figure}

\begin{figure}[H]
    \centering
    \includegraphics[width=0.95\textwidth]{include/SuttonExtended_S_E_Big.png}
    \includegraphics[width=0.95\textwidth]{include/SuttonExtended_S_U_Big.png}
    \caption{Comparing ADDQ and QL on two-sided bandit MDP with different variances on the left side and different exploration settings.}
    \label{fig:13}
\end{figure}

\newpage

\section{Grid world details}\label{app:gred}
\subsection{Experiment from the main text}
For reimplementation purposes we collect here all required information on the environment and the training for all plots.

Environment:
\begin{itemize}
    \item $\gamma=0.9$
    \item white low stochastic high average reward region: rewards $-0.05$, $+0.05$ with equal probabilities
    \item gray high stochastic low average reward region: rewards $-2.1$, $+2$ with equal probabilities
    \item goal: deterministic reward of $1$, fake goal: deterministic reward of $0.65$
\end{itemize}
Distributional properties for CategoricalQ, CategoricalDouble, and ADDQ:
\begin{itemize}
    \item categorical parametrization, 51 atoms equally spaced on $[-3,3]$
    \item initialization as $\delta_0$
\end{itemize}
Algorithmic choices:
\begin{itemize}
    \item $\beta$ is chosen according to \eqref{adap_rule}
    \item step-size schedule $\alpha_t(s,a)=\frac{1}{T_{s,a}(t)}$, with $T_{s,a}(t)$ the number of visits in $(s,a)$ up to time $t$, i.e. $\frac{1}{n}$ state-action wise counted,
    \item exploration: $\varepsilon$-greedy with $\varepsilon$ linearly decreasing from $1$ to $0.1$ in $10000$ steps, then constant
\end{itemize}
Policy evaluation:
     evaluation of 6 steps with a frequency of every 500 steps, correct action rates refer to if the exploration was greedy


\begin{figure}[h]
    \centering
    \includegraphics[width=0.95\linewidth]{include/Plot_2.1.png}
    \caption{additional plots}
\end{figure}
For completeness we give plots pairing learning curves for $Q$-values with the corresponding sample variance. We give all state-action combinations for the interesting states 1 and 4 next to the fake goal, 6 and 7 next to the region with high stochasticity, and 10 and 14 before leaving the region with high stochasticity.
\begin{figure}
    \centering
    \includegraphics[width=0.7\linewidth]{include/Plot_2.2.png}
     \includegraphics[width=0.7\linewidth]{include/Plot_2.3.png}
\end{figure}

\newpage

\begin{figure}[h]
    \centering    
    \label{fig:enter-label} \includegraphics[width=0.7\linewidth]{include/Plot_2.4.png}
\end{figure}

Finally, the following plot demonstrates that, in this GridWorld example, the relative variances are also strongly determined by the variances of the next state.

\begin{figure}[H]
    \centering
    \includegraphics[width=0.7\textwidth]{include/Relative_Variances.png}
\end{figure}

\subsection{Ablation study}\label{app:ablation}
In the following, on the same environment as before, we experiment with different choices for $\beta$ and compare with weighted DQL \cite{weightedDoubleQ}  with $c=10$ (WDQ), the standard choice from that paper. The choice of beta's names are comprised as follows:
\begin{itemize}
    \item (Optional) First two letters: Left-tilted (lt), Right-tilted (rt)
    \item First/Third letter: Neutral (n), Aggressive (a), Conservative (c)
    \item Final digit: Refers to the number of intervals in the definition of Beta (3 or 5)
\end{itemize}
The choices of Aggressive, Conservative, and Neutral refer to the trade-off of no interpolation (just choosing which Algorithm's update to take) and softening the interpolation, with neutral being in between the two choices and corresponding to the choice presented in the main body. Left- and Right-tilted refers to shifting the intervals for the relative Variance to fall into while choosing the interpolation coefficient. Left-tilted favors the Q update, Right-tilted the DQ update.\\
The concrete choices are:
\begin{align*}
    \text{n3:} \quad\quad 
    &\beta:= 
    \begin{cases}
    0.75&: S^2_{rel}(s,a) < 0.75 \\
    0.5&: S^2_{rel}(s,a)\in[0.75,1.25] \\
    0.25&: S^2_{rel}(s,a) > 1.25
    \end{cases}, \quad\quad\quad\quad
    &\text{a3:} \quad\quad
    &\beta:= 
    \begin{cases}
    1&: S^2_{rel}(s,a) < 0.99 \\
    0.5&: S^2_{rel}(s,a)\in[0.99,1.01] \\
    0&: S^2_{rel}(s,a) > 1.01
    \end{cases}\\
    \text{ltn3:} \quad\quad 
    &\beta:= 
    \begin{cases}
    0.75&: S^2_{rel}(s,a) < 1.25 \\
    0.5&: S^2_{rel}(s,a)\in[1.25,1.75] \\
    0.25&: S^2_{rel}(s,a) > 1.75
    \end{cases} , \quad\quad\quad\quad
    &\text{lta3:} \quad\quad 
    &\beta:= 
    \begin{cases}
    1&: S^2_{rel}(s,a) < 1.49 \\
    0.5&: S^2_{rel}(s,a)\in[1.49,1.51] \\
    0&: S^2_{rel}(s,a) > 1.51
    \end{cases}\\
    \text{rtn3:} \quad\quad 
    &\beta:= 
    \begin{cases}
    0.75&: S^2_{rel}(s,a) < 0.25 \\
    0.5&: S^2_{rel}(s,a)\in[0.25,0.75] \\
    0.25&: S^2_{rel}(s,a) > 0.75
    \end{cases}, \quad\quad\quad\quad
    &\text{rta3:} \quad\quad 
    &\beta:= 
    \begin{cases}
    1&: S^2_{rel}(s,a) < 0.49 \\
    0.5&: S^2_{rel}(s,a)\in[0.49,0.51] \\
    0&: S^2_{rel}(s,a) > 0.51
    \end{cases}\\
    \text{c3:} \quad\quad 
    &\beta:= 
    \begin{cases}
    0.6&: S^2_{rel}(s,a) < 0.6 \\
    0.5&: S^2_{rel}(s,a)\in[0.6,1.4] \\
    0.4&: S^2_{rel}(s,a) > 1.4
    \end{cases}, \quad\quad\quad\quad
    &\text{n5:} \quad\quad 
    &\beta:= 
    \begin{cases}
    1&: S^2_{rel}(s,a) \leq 0.25\\
    0.75&: S^2_{rel}(s,a)\in(0.25,0.75) \\
    0.5&: S^2_{rel}(s,a)\in[0.75,1.25] \\
    0.25&: S^2_{rel}(s,a)\in(1.25,1.75) \\
    0&: S^2_{rel}(s,a) \geq 1.75
    \end{cases}\\
    \text{ltc3:} \quad\quad 
    &\beta:= 
    \begin{cases}
    0.6&: S^2_{rel}(s,a) < 1.1 \\
    0.5&: S^2_{rel}(s,a)\in[1.1,1.9] \\
    0.4&: S^2_{rel}(s,a) > 1.9
    \end{cases}, \quad\quad\quad\quad
    &\text{ltn5:} \quad\quad 
    &\beta:= 
    \begin{cases}
    1&: S^2_{rel}(s,a) \leq 0.75\\
    0.75&: S^2_{rel}(s,a)\in(0.75,1.25) \\
    0.5&: S^2_{rel}(s,a)\in[1.25,1.75] \\
    0.25&: S^2_{rel}(s,a)\in(1.75,2.25) \\
    0&: S^2_{rel}(s,a) \geq 2.25
    \end{cases}\\
    \text{rtc3:} \quad\quad 
    &\beta:= 
    \begin{cases}
    0.6&: S^2_{rel}(s,a) < 0.1 \\
    0.5&: S^2_{rel}(s,a)\in[0.1,0.9] \\
    0.4&: S^2_{rel}(s,a) > 0.9
    \end{cases},\quad\quad\quad\quad
    &\text{rtn5:} \quad\quad 
    &\beta:= 
    \begin{cases}
    1&: S^2_{rel}(s,a) \leq -0.25\\
    0.75&: S^2_{rel}(s,a)\in(-0.25,0.25) \\
    0.5&: S^2_{rel}(s,a)\in[0.25,0.75] \\
    0.25&: S^2_{rel}(s,a)\in(0.75,1.25) \\
    0&: S^2_{rel}(s,a) \geq 1.25
    \end{cases}
\end{align*}
\begin{align*}
    \text{a5:} \quad\quad 
    &\beta:= 
    \begin{cases}
    1&: S^2_{rel}(s,a) \leq 0.99\\
    0.75&: S^2_{rel}(s,a)\in(0.99,0.995) \\
    0.5&: S^2_{rel}(s,a)\in[0.995,1.005] \\
    0.25&: S^2_{rel}(s,a)\in(1.005,1.01) \\
    0&: S^2_{rel}(s,a) \geq 1.01
    \end{cases}, \quad\quad\quad\quad
    &\text{c5:} \quad\quad 
    &\beta:= 
    \begin{cases}
    0.7&: S^2_{rel}(s,a) \leq 0.1\\
    0.6&: S^2_{rel}(s,a)\in(0.1,0.7) \\
    0.5&: S^2_{rel}(s,a)\in[0.7,1.3] \\
    0.4&: S^2_{rel}(s,a)\in(1.3,1.9) \\
    0.3&: S^2_{rel}(s,a) \geq 1.9
    \end{cases}\\
    \text{lta5:} \quad\quad 
    &\beta:= 
    \begin{cases}
    1&: S^2_{rel}(s,a) \leq 1.49\\
    0.75&: S^2_{rel}(s,a)\in(1.49,1.495) \\
    0.5&: S^2_{rel}(s,a)\in[1.495,1.505] \\
    0.25&: S^2_{rel}(s,a)\in(1.505,1.51) \\
    0&: S^2_{rel}(s,a) \geq 1.51
    \end{cases}, \quad\quad\quad\quad
    &\text{ltc5:} \quad\quad 
    &\beta:= 
    \begin{cases}
    0.7&: S^2_{rel}(s,a) \leq 0.6\\
    0.6&: S^2_{rel}(s,a)\in(0.6,1.2) \\
    0.5&: S^2_{rel}(s,a)\in[1.2,1.8] \\
    0.4&: S^2_{rel}(s,a)\in(1.8,2.4) \\
    0.3&: S^2_{rel}(s,a) \geq 2.4
    \end{cases}\\
    \text{rta5:} \quad\quad 
    &\beta:= 
    \begin{cases}
    1&: S^2_{rel}(s,a) \leq 0.49\\
    0.75&: S^2_{rel}(s,a)\in(0.49,0.495) \\
    0.5&: S^2_{rel}(s,a)\in[0.495,0.505] \\
    0.25&: S^2_{rel}(s,a)\in(0.505,0.51) \\
    0&: S^2_{rel}(s,a) \geq 0.51
    \end{cases}, \quad\quad\quad\quad
    &\text{rtc5:} \quad\quad 
    &\beta:= 
    \begin{cases}
    0.7&: S^2_{rel}(s,a) \leq -0.4\\
    0.6&: S^2_{rel}(s,a)\in(-0.4,0.2) \\
    0.5&: S^2_{rel}(s,a)\in[0.2,0.8] \\
    0.4&: S^2_{rel}(s,a)\in(0.8,1.4) \\
    0.3&: S^2_{rel}(s,a) \geq 1.4
    \end{cases}
\end{align*}
\newpage
\begin{figure}[H]
    \centering
    \includegraphics[width=\textwidth]{include/Ablation_study_1_Biases.png}
    \label{fig:1}
\end{figure}
\begin{figure}[H]
    \centering
    \includegraphics[width=\textwidth]{include/Ablation_study_2_Biases.png}
    \caption{Ablation study plots regarding bias improvement. State 1 is adjacent to the fake goal, state 6 adjacent to the stochastic region.}
    \label{fig:2}
\end{figure}
It turns out that
\begin{itemize}
    \item the choice of thresholds in $\beta$ (hyperparameter to the algorithm) is harmless, results do not vary a lot,
    \item conservative choices seem to work especially well.
    \item weighted double Q-learning, an algorithm that is similar to ADDQ with non-distributional choice of $\beta$ is improved by our locally adaptive distributional RL based choice of $\beta$.
\end{itemize}

\subsection{Comparison to other algorithms}\label{sec:comparison}

We compared ADDQ with other bias reduction methods using different hyperparameters used in the respective papers on the same GridWorld environment.
\begin{itemize}
    \item Maxmin with 2, 4, 6, and 8 ensembles \cite{Maxmin}
    \item Ensemble Bootstrapped QL (EBQL) with 3, 7, 10, and 15 ensembles \cite{EnsembleQ}
    \item Randomized Ensemble DQL (REDQ) with 3 and 5 ensembles and size 1 and 2 of random update subset sizes \cite{redq}
\end{itemize}
It turns out that ADDQ decreases the estimation bias stronger than those algorithms while only using two ensembles.

\clearpage

\begin{figure}[H]
    \centering
    \includegraphics[width=\textwidth]{include/Comparison_to_MaxMin_Scores.png}
    \includegraphics[width=\textwidth]{include/Comparison_to_MaxMin_Q_Values.png}
    \caption{Comparison to MaxMin.}
\end{figure}

\begin{figure}[H]
    \centering
    \includegraphics[width=\textwidth]{include/Comparison_to_EBQL_Scores.png}
    \includegraphics[width=\textwidth]{include/Comparison_to_EBQL_Q_Values.png}
    \caption{Comparison to EBQL}
\end{figure}

\begin{figure}[H]
    \centering
    \includegraphics[width=\textwidth]{include/Comparison_to_REDQ_Scores.png}
    \includegraphics[width=\textwidth]{include/Comparison_to_REDQ_Q_Values.png}
    \caption{Comparison to REDQ.}
\end{figure}

\section{ADDQ adaptation for QRDQN - setup}\label{appendix:QRDQN}
To not double too much in the main text we moved the adaptation of ADDQ to the quantile setup \cite{dabney2018} to the appendix. In this section we explain how to adapt the ADDQ idea into QRDQN, experimental results are provided in the next section. 

The categorical approach has multiple disadvantages, most notably rewards and the fixed atom positions must be compatible. The categorical algorithm was included in the main text to keep notation simple. For quantile distributional RL the return distributions are parametrized by
\[\mathcal{F}_{Q,m} = \Big\{  \sum_{i=1}^{m}\frac{1}{m} \delta_{\theta_i} : \theta_i \in \mathbb{R} \Big\}.\]
In contrast to the categorical setting the positions of the atoms are not fixed, but the weights of all atoms are set equal. The computation of the target is equal to the categorical version (except using the projection on $\mathcal F_{Q,m}$ instead of $\mathcal F_{C,m}$). The update step is a gradient step in computing the Wasserstein-projection on $\mathcal F_{Q,m}$ of the target distribution $\hat\eta$, that is a gradient step in the quantile Huber-loss minimization:
\begin{align*}
\min_{\hat{\theta}_{1}^A(s,a),...,\hat{\theta}_{m}^A(s,a)} \sum_{i=1}^{m}\mathbb{E}_{Z \sim \hat{\eta}}[\rho_{\tau_i}^{\kappa}(Z-\hat{\theta}_{i}^A(s,a))],
\end{align*}
with quantile mid-points $\tau_i=\frac{2i-1}{2m}$ and
\begin{align*}
    \rho_{\tau}^{\kappa}(u) 
    =
    \begin{cases} 
    |\tau - \textbf{1}_{u<0}|\frac{1}{2}u^2&:|u|\leq \kappa \\
     |\tau - \textbf{1}_{u<0}|\kappa(|u|-\frac{1}{2}\kappa) &: |u|> \kappa
     \end{cases}.
\end{align*}
The quantile $Q$-learning modification of classical DQN \cite{dqn} is called QRDQN. Using the same network architecture as C51, QRDQN approximates return distributions using the quantile representation. Therefore the last layer outputs the $m$ quantile locations for each action.
In the quantile setup we write $\eta_\omega(s,a)=\frac{1}{m}\sum_{i=1}^m\delta_{\theta_i(s,a;\omega)}$ with induced mean values $Q_\omega(s,a)=\frac{1}{m}\sum_{i=1}^m\theta_i(s,a;\omega)$. Given a sample transition $(s,a,r,s')$, the network parameters are updated via gradient descent with respect to the loss function
\begin{align}\label{eq:qrdqn loss}
    \mathcal L(\omega)&=\frac{1}{m}\sum_{i,j=1}^m\rho_{\tau_i}^1(r+\gamma\theta_j(s',z^*;\bar\omega)-\theta_i(s,a;\omega)), \quad
    z^*=\operatorname{argmax}_{a'}Q_{\bar\omega}(s',a'),
\end{align}
and the quantile mid-points $\tau_i=\frac{2-1}{2m}$. 

In what follows we turn three known double variants of DQN with overestimation reduction into QRDQN variants and compare them on several Arcade environments to our algorithms. We use double DQN \cite{hasselt2015doubledqn}, a $Q$-learning adaptation of the clipping trick included in TD3 \cite{fujimoto2018addressing}, a quantile version of our ADDQ algorithm, and an additional variant of ADDQ. Using the Stable-Baselines3 framework [\cite{sb3}] there is very little that must be modified. Return distributions $\eta^A$ and $\eta^B$ are parametrized  with two independently initialized neural networks denoted by $\omega^A$ and $\omega^B$. As in DQN we use delayed target networks, one for $A$, one for $B$, that are indicated with an additional bar. For each gradient step we simulate a vector of random variables with the same size as the batch size with each element determining which of the two estimators is being updated based on the respective transition with the same position in the batch. Accordingly, we use twice the batch size for these methods, so that on average per gradient step, the same number of transitions is used for each estimator, compared to the single-estimator case. The only difference in different algorithms is the target used to update the neural networks.

Similar to modifying C51 implementations we only modify the target return distributions $b_{r,\gamma}\#\eta_{\bar\omega}(s',z^*)$ for an appropriate action $z^*$. In the quantile setup those are given by the locations of their atoms:
\begin{align*}
    \Gamma=\{r+\gamma\theta_j(s',z^*;\omega):j=1,\dots,m\}    
\end{align*}
We again use the compact $A/B$ notation to indicate how the update applies for $\Gamma^A$ and $\Gamma^B$.


\textbf{Double QRDQN:} $\Gamma^{A/B} = \eta_{\bar\omega^{B/A}}(s',z^*)$, where  $z^*=\operatorname{argmax}_{a'}Q_{\bar\omega^{A/B}}(s',a')$. \smallskip

%\textbf{(ii) Double estimator based on target network}, inspired by \cite{hasselt2015doubledqn}: The batch size is the same as in the single estimator case, there are no extra networks. The choice corresponds to 
 %   \begin{align*}
  %      \Gamma &= \eta_{\bar\omega}(s',z^*),\quad z^*=\operatorname{argmax}_{a'}Q_{\omega}(s',a').
   % \end{align*}
    %Greedy action selection is with respect to the online network, i.e. the target network acts as B, the online network as A.\smallskip

   \textbf{Clipped QRDQN:}  $\Gamma^{A/B}=\eta_{\bar\omega^X}(s',z^*)$, where
   $z^*=\operatorname{argmax}_{a'}Q_{\bar\omega^{A/B}}(s',a')$ and $X=\operatorname{argmin}_{c\in\{A,B\}}Q_{\bar\omega^c}(s',z^*)$.  \smallskip


%\textbf{Truncated QRDQN}, inspired by the algorithm TQC \cite{tqc}: Atoms of both target distributions are combined, the top $k$ atoms are dropped to introduce a negative bias:
 %   \begin{align*}
  %      \Gamma^{A/B}&=\{r+\gamma\tilde\theta_j(s',z^*):j=1,\dots,(2m-k)\},
   % \end{align*}
    %where
    %\begin{align*}        
     %   \tilde\theta_j(s',a')&=\text{sort}(\{\theta_i(s,a;\bar\omega^X): i=1,\dots,m,\ X\in\{A,B\}\})_j
    %\end{align*}
%and $z^*=\operatorname{argmax}_{a'}Q_{\bar\omega^{A/B}}(s',a')$ and accordingly the $m$ in (\ref{eq:qrdqn loss}) is replaced by $2m-k.$\smallskip


    \textbf{ADDQ (us):} 
$\Gamma^{A/B}=\beta\eta_{\bar\omega^{A/B}}(s',z^*)+(1-\beta)\eta_{\bar\omega^{B/A}}(s',z^*)$,
where $z^*=\operatorname{argmax}_{a'}Q_{\bar\omega^{A/B}}(s',a').$ The weights $\beta=\beta(s,a;\omega)$ are essentially arbitrary and can depend on the current estimated return distributions $\eta^A$ and $\eta^B$. For the experiments we take the same choice from \ref{adap_rule} as used for the tabular and the categorical settings.

%To show robustness of our adaptive bias control approach we introduce a second (more indirect) way for adaptive bias control. We need this alternative for continuous control as $\beta$ from Equation \eqref{adap_rule} is not well-defined for infinitely many actions.

%\textbf{(v) Wasserstein reduced mixture (Wasserstein QRDQN, us):}  In order to control the overestimation bias, based on the current estimates, we use the 1-Wasserstein $w_1$ distance of both estimates from the mixture distribution. A large discrepancy between $\eta_{\bar\omega^A}$ and $\eta_{\bar\omega^B}$ is a sign that one estimate has been overestimated. Since in DRL we have access to the full return distribution, this allows us to measure the estimator's divergence directly based on the Wasserstein distance, a metric for probability distributions. The more aligned both distributions, the smaller the reduction. We set
%\begin{align*}
%\Gamma=\{r+\gamma(\theta_j(s',z^*;\bar\omega^X)-\beta w_1(\eta_{\bar\omega^A}(s',z^*),\eta_{\bar\omega^B}(s',z^*))):j=1,\dots,m,\ X=A,B\},
%\end{align*}
%where $z^*=\operatorname{argmax}_{a'}(Q_{\bar\omega^A}(s',a')+Q_{\bar\omega^B}(s',a'))$ and accordingly the $m$ in (\ref{eq:qrdqn loss}) is replaced by $2m.$ Since $w_1(\nu,\nu')\geq|\mathbb E_{Z\sim\nu}[Z]-\mathbb E_{Z\sim\nu'}[Z]|$ factors fulfilling $\beta\in(0,\frac{1}{2})$ are reasonable. In the Atari experiments we fix $\beta=0.1$.


%Here, for the double estimator variants, the target return distribution $\Gamma = \eta_{\bar\omega}(s',a^*)=\frac{1}{m}\sum_{j=1}^m\delta_{\theta_j(s',a^*;\bar\omega)}$ is replaced as follows:
%\begin{itemize}
%    \item double, clipped, LTV adaptive, analogously to the above,
%    \item Wasserstein reduced mixture:
%    \[\Gamma=\frac{1}{2m}\sum_{j\in\{1,\dots,m\},X\in\{A,B\}}\delta_{\theta_{j}(s',z^*;\bar\omega^X)-\beta w_1(\eta_{\bar\omega^A}(s',z*),\eta_{\bar\omega^B}(s',z*))},\]
%    where $z^*=\argmax_{a'}Q_{\bar\omega^{A/B}}(s',a').$
%    \item truncated mixture, dropping the topmost $k$ atoms of the mixture: 
%    \begin{align*}
%        \Gamma^{A/B}&=\frac{1}{2m-k}\sum_{j=1}^{2m-k}\delta_{\hat\theta_j(s',z^*)} \quad \text{where}\\
%        \hat\theta_j(s',a')&=\text{sort}(\{\theta_i(s,a;\bar\omega^X)| i=1,\dots,m,\ X\in\{A,B\}\})_j\\
%    \end{align*}
%    and $z^*=\argmax_{a'}Q_{\bar\omega^{A/B}}(s',a').$
%\end{itemize}

%\begin{figure}[h]
%	\centering
%	\ \includesvg[width=0.95\textwidth]{qrdqn}
%	\caption{QRDQN platzhalter, pro environment 7 plots, falls lesbar (C51, double, double 2015, truncated, clip, ltv, ws-reduced),das wird noch besser aussehen hoffentlich, hier teilweise nur 1 Seed erstmal}
%	\label{fig:QRDQN}
%\end{figure}
Experimental results are presented in the next section.

\newpage

\section{Deep reinforcement learning experiments}\label{app:experiments}

To ensure fair comparison we modified the algorithms C51 [\cite{bellemare2017}] and QRDQN [\cite{dabney2018}] within the Stable-Baselines3 framework [\cite{sb3}]. The C51 implementation has been added to this framework by adapting from the Dopamine framework [\cite{dopamine}] and the DQN Zoo [\cite{dqnzoo}]. We run Atari environments from the Arcade Learning Environment [\cite{ale}] and MuJoCo [\cite{mujuco}] environments both using the Gymnasium API [\cite{gymnasium}]. We run the experiments via the RL Baselines3 Zoo [\cite{rl-zoo3}] training framework.

The experiments were executed on a HPC cluster with NVIDIA Tesla V100 and NVIDIA A100 GPUs. The replay buffer on Atari environments takes around 57GB of memory and less than 7 GB of memory for MuJoCo environments.

For the experiments the training has been interrupted every 50000 steps and 10 evaluation episodes on 10 evaluation environments without exploration have been performed. The plots below show the mean total reward (sum of all rewards) averaged over \textbf{10 seeds} with standard errors of seeds as the shaded regions. To improve visibility a rolling window of size 4 is applied. Atari runs took less than 48 hours for 20 million train steps and periodic evaluations and MuJoCo runs less than 36 for 10 million train steps (Humanoid) and periodic evaluations. Note that one timestep in the Atari environments corresponds to 4 frames, which are stacked together. This corresponds to repeating every action 4 times in the actual game. Therefore 20 million timesteps correspond to 80 million frames. Additionally, a small ablation study comprising of 10 seeds on one evaluation environment with some of the choices of beta detailed in Appendix \ref{app:ablation} has been conducted.

\newpage
\subsection{Full experimental results for the categorical parametrization}
As in \cite{bellemare2017} we use 51 atoms for all C51 variants.
\begin{figure*}[h]
\begin{center}
\includegraphics[height=8cm]{include/C51one_fig_all_graphs.png}
\includegraphics[height=8cm]{include/C51one_fig_all_graphs_2.png}
\caption{Learning curves on 10 Atari environments, averaged over 10 seeds}
\end{center}
\end{figure*}
\newpage
We additionally provide plots using the RLiable library \cite{RLiable}.
\begin{figure*}[h]
\begin{center}
 \includegraphics[width=0.6\textwidth]{include/C51_Probability_of_Improvement.png}
        \caption{RLiable probability of improvement plot}
    \vspace{5mm}        
    \includegraphics[width=0.95\textwidth]{include/C51Rliable_metrics.png}
    \caption{RLiable human normalized scores plot (based on \cite{agent57})}
    \end{center}
\end{figure*}


\newpage
\subsection{Full experimental results for the quantile parametrization}
\begin{figure*}[h]
\begin{center}
\includegraphics[height=8cm]{include/QRone_fig_all_graphs_1.png}
\includegraphics[height=8cm]{include/QRone_fig_all_graphs_2.png}
\caption{Learning curves on 10 Atari environments, averaged over 10 seeds}
\end{center}
\end{figure*}
\newpage

\begin{figure*}[h]
\begin{center}
     \includegraphics[width=0.6\textwidth]{include/QRProbability_of_Improvement.png}
        \caption{RLiable probability of improvement plot}
    \vspace{5mm}
        
    \includegraphics[width=0.95\textwidth]{include/QRRliable_metrics.png}
    \caption{RLiable human normalized scores plot (based on \cite{agent57})}
  
    \end{center}
\end{figure*}

\newpage 


\subsection{Full experimental results for the actor critic setting}

\begin{figure*}[h]
\begin{center}
\includegraphics[width=0.9\textwidth]{include/one_fig_all_graphs.png}
\caption{Learning curves on 5 MuJoCo environments, averaged over 10 seeds}
\vspace{5mm}
\includegraphics[width=0.5\textwidth]{include/SACProbability_of_Improvement.png}
        \caption{RLiable probability of improvement plot}
    \end{center}
\end{figure*}

\begin{figure*}[h]
\begin{center}
\includegraphics[width=0.7\textwidth]{include/SACRliable_metrics.png}
    \caption{RLiable normalized scores plot,  based on highest and lowest performance on each environment}
\end{center}
\end{figure*}

\newpage 
\quad
\newpage



%%%%%%%%%%%%%%%%%%%%%%%%%%%%%%%%%%%%%%%%%%%%%%%%%%%%%%%%%%%%%%%%%%%%%%%%%%%%%%%
%%%%%%%%%%%%%%%%%%%%%%%%%%%%%%%%%%%%%%%%%%%%%%%%%%%%%%%%%%%%%%%%%%%%%%%%%%%%%%%



\section{Convergence proof in the tabular categorical setup}\label{sec:proof}
In this section we give a convergence proof for the adaptive distributional double-Q algorithm in the simplest setting, the categorical setting. The proof is based on known arguments from the literature and requires some modifications to work in our generality. Since many papers only sketched proofs we decided to spell out all details.
\begin{remark}[Notation and short recap]
The Cramer distance $\ell_2$ for probability distributions $\nu,\nu'\in\cP(\R)$ is given by
\[\ell_2(\nu, \nu') = \left( \int_{\mathbb{R}} \left| F_{\nu}(z) - F_{\nu'}(z) \right|^2 dz \right)^{1/2}.\]
Following \cite{rowland18,drl} the supremum extension of a probability metric $d$ between two return distribution functions $\eta,\eta'\in\cP^{\cS\times\cA}$ is denoted as
\[\bar d(\eta,\eta')=\sup_{s,a\in\cS\times\cA}d(\eta(s,a),\eta'(s,a).\]
The iterates $\eta_{k+1}=\Pi_C\cT^\pi\eta_k$ converge to the unique fixed point in $\cF_{C,m}^{\cS\times\cA}$ with respect to $\bar\ell_2$ based on Banach's fixed point Theorem. This follows from the contraction property
\begin{equation}
\bar\ell_2(\Pi_C\cT^\pi\eta,\Pi_C\cT^\pi\eta')\leq\sqrt{\gamma}\bar\ell_2(\eta,\eta'),
\end{equation}
[compare \cite{rowland18,drl}].
\end{remark}

\begin{theorem}[Convergence of adaptive distributional $Q$-learning in the categorical setting]\label{thm:double categorical control}
    Given some initial return distribution functions $\eta_0^{A},\eta^{B}_0$ supported within $[\theta_1,\theta_m]$, the induced $Q$-values, i.e. the expected values of the return distributions $(\eta_t^{A}),(\eta^{B}_t)$, recursively defined by Algorithm \ref{alg:CDRL} converge almost surely towards $Q^*$ if the following conditions are satisfied:
\begin{enumerate}
	\item the step sizes $\alpha_t(s,a)$ almost surely fulfill the Robbins-Monro conditions $\sum_{t=0}^{\infty}\alpha_t(s,a)=\infty$ and $
		\sum_{t=0}^{\infty}\alpha_t^{2}(s,a)<\infty$.
	\item rewards are bounded in $[R_{min},R_{max}]$ and $[\frac{R_{min}}{1-\gamma},\frac{R_{max}}{1-\gamma}]\subseteq[\theta_1,\theta_m],$
    \item the choice of updating $\eta^A$ or $\eta^B$ is random and independent of all  previous random variables 
    \item the sequences $(\beta^A_t)_{t\in\N},(\beta^B_t)_{t\in\N}$ only depend on the past and fulfill $\lim_{t\to\infty} |\beta^A_t-\beta^B_t|=0$ almost surely.
\end{enumerate}
If additionally the MDP has a unique optimal policy $\pi^*,$ then $(\eta_t^{A}),(\eta^{B}_t)$ converge almost surely in $\bar\ell_2$ to some limit $\eta^*_C\in\cF_{C,m}$ and the greedy policy with respect to $\eta^{*}_C$ is the optimal policy.
\end{theorem}

Note that the algorithm and proof uses $\beta^{A/B}_{t+1}(s,a)$ with index ${t+1}$ when updating $\eta^{A/B}_t$. This is to show that in general the parameter is allowed to depend on $S_{t+1}$ and the respective greedy action, i.e. it must only be $\cF_{t+1}$ measurable. To portray this generality in the following we will only write $\beta^{A/B}_{t+1}$ without referencing a state-action pair.

The simplest way to guarantee the assumptions on the adaptive parameters $\beta^A$, $\beta^B$ to be satisfied is to chose them equal.\smallskip

As in \cite{hasselt2010}, the proof is based on the following stochastic approximation result, see also  \cite{bertsekasbook}, Proposition 4.5. %%%TODO: SARSA Quelle (van Hasselt2015 double hat keine Info zu Journal in dem es veröffentlich ist, sonst glaube ich keine Fehler zu bibtex Dingen)
\begin{lemma}[\cite{singh2000}, Lemma 1]\label{lem:stochastic approximation lemma} %übernommen von Leo's Masterarbeit
Suppose \((\Omega, \mathcal{A}, \mathbb{P}, (\mathcal{F}_n))\) is a filtered probability space on which all appearing random variables are defined. Suppose that

\begin{enumerate}
    \item a stochastic process \((F_n)_{n \in \mathbb{N}} \subset \mathbb{R}^d\) with the coordinates \(F_{i,n}\) for \(i = 1, \ldots, d\) such that \(F_n\) is \(\mathcal{F}_{n+1}\)-measurable and for all \(i = 1, \ldots, d\)
    \[
    \left\|\mathbb{E}[F_n | \mathcal{F}_n]\right\|_\infty \leq \kappa\|X_n\|_\infty + c_n \quad \text{and} \quad \mathbb{V}[F_{i,n} | \mathcal{F}_n] \leq K(1 + \kappa\|X_n\|_\infty)^2 \quad n \geq 1,
    \]
    where \(\kappa \in [0,1)\), an adapted, stochastic process \((c_n)_{n \in \mathbb{N}} \subset \mathbb{R}^+\) that converges to 0 almost surely and some constant \(K > 0\).

    \item the non-negative stochastic process \((\alpha_n)_{n \in \mathbb{N}} \subset \mathbb{R}^d\), with the coordinates \(\alpha_{i,n} \in [0,1]\) for \(i = 1, \ldots, d\) is adapted with
    \[
    \sum_{n=1}^{\infty} \alpha_{i,n} = \infty \quad \text{and} \quad \sum_{n=1}^{\infty} \alpha_{i,n}^2 < \infty \quad \text{a.s.}.
    \]
\end{enumerate}

Then, for any \(\mathcal{F}_0\)-measurable initial condition \(X_0\) the stochastic process \((X_n)_{n \in \mathbb{N}} \subset \mathbb{R}^d\) with coordinates \(X_{i,n}\) for \(i = 1, \ldots, d\) that is recursively defined by
\[
X_{i,n+1} = (1 - \alpha_{i,n})X_{i,n} + \alpha_{i,n}F_{i,n}, \quad n \in \mathbb{N},
\]
converges almost surely to zero.
\end{lemma}
Furthermore, we follow \cite{rowland18} by first showing the convergence of the mean-values to $Q^*$ and afterwards showing convergence of the return distribution functions, under the assumption of a unique optimal policy, by coupling it with policy evaluation. The convergence of the latter is easier to prove and we will do so at the end.
\begin{lemma}[Adaptive Double Categorical Temporal Difference for Policy Evaluation]\label{lem: double categorical policy evaluation}
    Given some initial return distribution functions $\eta_0^{A},\eta^{B}_0$ supported within $[\theta_1,\theta_m]$ and a stationary policy $\pi\in\Pi_S$, the return distribution functions $(\eta_t^{A}),(\eta^{B}_t)$ recursively defined by Algorithm \ref{alg:CDRL}, but with $a^*\sim\pi(\cdot;S_{t+1})$ instead,
    converge almost surely towards the unique fixed point $\eta_C\in\cP(\R)^{\cS\times\cA}$ of the operator $\Pi_C\cT^\pi$ with respect to $\bar\ell_2$, if the following conditions are satisfied:
\begin{enumerate}
	\item the step sizes $\alpha_t(s,a)$ fulfill the Robbins-Monro conditions:
	\begin{itemize}
		\item[$\bullet$] $\sum_{t=0}^{\infty}\alpha_t(s,a)=\infty$
		\item[$\bullet$]$\sum_{t=0}^{\infty}\alpha_t^{2}(s,a)<\infty,$
	\end{itemize}
	\item rewards are bounded in $[R_{min},R_{max}]$ and $[\frac{R_{min}}{1-\gamma},\frac{R_{max}}{1-\gamma}]\subseteq[\theta_1,\theta_m],$
    \item the choice of updating $\eta^A$ or $\eta^B$ is random and independent of all other previous random variables 
\end{enumerate}
\end{lemma}
The above result is only relevant for the proof of Theorem \ref{thm:double categorical control}, as policy evaluation with a double estimator is not of interest. Note that convergence of categorical temporal difference for policy evaluation (in the single estimator case) has been proven in [\cite{rowland18} Theorem 2 mimicking \cite{Tsitsiklis1994} Theorem 2] and [\cite{drl} Theorem 6.12 applying \cite{Tsitsiklis1994} Theorem 3 or \cite{bertsekasbook} Proposition 4.5]. 
\begin{lemma}\label{lem:RM}
    Let $(\alpha_t)_{t\in\N_0}$ be a sequence fulfilling the Robbins-Monro conditions and $(Y_t)_{t\in\N}$ an $iid$ sequence of Bernoulli(0.5) random variables, i.e. $\mathbb P(Y_t=1)=\mathbb P(Y_t=0)=0.5$ for all $t\in\N_0$.
    Then $(\alpha_tY_t)_{t\in\N_0}$ also fulfills the Robbins-Monro condition.
\end{lemma}
\begin{proof}
    The almost sure convergence of the summed squares is obviously fulfilled due to 
    \[\sum_{t=0}^\infty(\alpha_tY_t)^2\leq\sum_{t=0}^\infty \alpha_t^2<\infty\quad \text{almost surely.}\]
    Due to independence of each $Y_t$ with $\{Y_n|n\in\N_0,\ n\neq t\}$ as well as with $\alpha=(\alpha_t)_{t=0}^\infty$ we will consider a two stage experiment, where we first draw the sequence $\alpha=(\alpha_t)_{t=0}^\infty$ and then independently of this realization sample the $iid$ sequence $Y=(Y_t)_{t=0}^\infty$. Due to the independence the joint measure of $\alpha$ and $Y$ is the product measure. Consider the product space $(\Omega,\cF,\mathbb P)=(\Omega_\alpha\times\Omega_Y,\cF_\alpha\otimes\cF_Y,\mathbb P^{\otimes\N}_\alpha\otimes\mathbb P^{\otimes\N}_Y)$ where $\Omega_\alpha,\Omega_Y=[0,1]^\N$, $\cF_\alpha,\cF_Y=\cB([0,1])^{\otimes \N}$ . Then, using that $\sum_{t=0}^\infty\alpha_t=\infty$ $\mathbb P_\alpha$-almost surely, we have
    \begin{align*}   
    \mathbb P\Big(\sum_{t=0}^\infty\alpha_tY_t=\infty\Big)
    =&\int_{\Omega_\alpha}
    \mathbb P_Y\Big(\sum_{t=0}^\infty\alpha_tY_t=\infty\Big)
    d\mathbb P_\alpha(\alpha)\\
            =&\int_{\{(\alpha_t)_{t=0}^\infty\in \Omega_\alpha:\sum_{t=0}^\infty\alpha_t=\infty\}}\mathbb P_Y\Big(\sum_{t=0}^\infty\alpha_tY_t=\infty\Big)d\mathbb P_\alpha(\alpha)\\
            \overset{(a)}{=}&\int_{\{(\alpha_t)_{t=0}^\infty\in \Omega_\alpha:\sum_{t=0}^\infty\alpha_t=\infty\}}1\ d\mathbb P_\alpha(\alpha)\\
            =& 1,
    \end{align*} 
    where $(a)$ can be seen as follows. Consider any deterministic sequence $(b_t)\subseteq[0,1]$ fulfilling $\sum_{t=0}^\infty b_t=\infty$. Then
    \[\infty= \sum_{t=0}^\infty b_t=\sum_{t=0}^\infty b_tY_t+\sum_{t=0}^\infty b_t\textbf{1}_{Y_t=0}.\]
    Now notice that $A=\sum_{t=0}^\infty b_tY_t$ and $B=\sum_{t=0}^\infty b_t\textbf{1}_{Y_t=0}$ are identically distributed and since the sum of $A$ and $B$ is always infinity, almost surely either one of them is infinite. Given the identical distribution, we infer
    \[\mathbb P_Y(\sum_{t=0}^\infty b_tY_t=\infty)>0.\]
    But since $(b_tY_t)$ is an independent sequence of random variables and the event that the infinite sum diverges is in the tail sigma algebra, the Kolmogorov 0-1 law yields:
    \[\mathbb P_Y(\sum_{t=0}^\infty b_tY_t=\infty)=1.\]
    %Now, for the first condition we recursively construct a deterministic sequence as follows. For $K\in\N$, define
    %\[t_K:=\min\{\n\in\N | n>t_{K-1}\text{ and }\sum_{t=0}^\infty\alpha_t>K \text{ almost surely}\}.\]
    %Since $\sum_{t=0}^\infty\alpha_t=\infty $ almost surely, such a sequence exists. Further define a deterministic sequence $(b_t)$ that is one exactly at the %indices $t_K$, for $K\in\N$, and $0$ otherwise. This is
    %\[b_t=\begin{cases}
    %    1\quad ,\ if\ \exists K\in\N,\ s.t.\ t=t_K\\
    %    0\quad otherwise.
    %\end{cases}\]
    %Then, we have
    %\[\sum_{t=0}^\infty\alpha_tY_t \]
\end{proof}
\begin{remark}
    As outlined in \cite{rowland18}, proof of Proposition 1, denoting by $\cM(\R)$ the space of all finite signed measures on $(\R,\cB(\R)),$ the subspace 
    \[\cM_0(\R):=\{\nu\in\cM(\R)|\nu(\R)=0,\ \int_\R F_\nu (x)^2dx<\infty \},\]
    "where $F_\nu (x)=\nu([-\infty,x))$ for $x\in\R$, is isometrically isomorphic to a subspace of the Hilbert space $L^2(\R)$ with inner product given by 
    \[\langle\nu_1,\nu_2\rangle_{\ell_2}=\int_\R F_{\nu_1}(x)F_{\nu_2}(x)dx."\]
    Then the affine translation $\delta_0+\cM_0$ is also Hilbert space endowed with the same inner product. It contains probability measures $\nu\in\cP(\R)$ satisfying
    \[\int_{-\infty}^{0} F_{\nu}(x)^2 \, dx < \infty \quad \text{and} \quad \int_{0}^{\infty} (1 - F_{\nu}(x))^2 \, dx < \infty.\]
    To see this, consider $\mu= \nu - \delta_0$ fulfills $F_\mu(x)=F_\nu(x)$ for $x<0$ and $F_\mu(x)=F_\nu(x)-1$ for $x\geq 1$. Hence, $\mu\in\cM_0.$ The two conditions assure that the tails decay fast enough.\\
    Note that the inner product induces a norm through $\|\nu\|^2_{\ell_2}=\langle\nu,\nu\rangle.$ And we have $\ell_2(\nu_1,\nu_2)=\|\nu_1-\nu_2\|_{\ell_2}$. In the following proof, we will make use of the relationship
    \[\ell_2^2(\nu_1+\nu_2,\nu'_1+\nu'_2)=\|\nu_1-\nu'_1\|^2_{\ell_2}+\|\nu_2-\nu'_2\|^2_{\ell_2}+2\langle\nu_1-\nu'_1,\nu_2-\nu'_2\rangle\]
    holding by bilinearity of the inner product.
\end{remark}
\begin{proof}[Proof of Thoerem \ref{thm:double categorical control}]
\textbf{Step 1: Convergence of mean values to $Q^*$}\\
The proof mainly follows \cite{rowland18} and \cite{hasselt2010}. Let the filtration be given by $\cF_t=\sigma(\eta^A_0,\eta^B_0,s_0,a_0,\alpha_0,R_0,S_1,Y_1,\beta^A_1,\beta^B_1\dots,s_t,a_t,\alpha_t),$ where $(Y_{n})_{n\in\N}$ is an iid sequence of \textit{Bernoulli(0.5)} random variables, independent of all other appearing random variables, such that A is updated when $Y_{n+1}=1$. Denote the expected values of the return-distributions by $Q^{A}_t(s,a)=\E_{R\sim\eta_t^A(s,a)}[R]$ and overloading notation, let us further write $\E[\nu]$ for the expected value $\E_{R\sim\nu}[R]$ of a probability distribution $\nu\in\cP(\R)$. We will first consider how the expected values evolve. Due to the symmetry of the updates it is sufficient to show convergence of $Q^A_t$ to $Q^*$. It is implied that $\alpha(s,a)=0$ for $(s,a)\neq(s_t,a_t)$.
Further, define
    \begin{align*}
        X_t(s_t,a_t) &:= Q^A_t(s_t,a_t)-Q^*(s_t,a_t)\\
        F_t(s_t,a_t) &:= \textbf{1}_{Y_{t+1}=1}\Big(R_t + \gamma(\beta^A_{t+1}Q_t^A(S_{t+1},a^*)+(1-\beta^A_{t+1})Q_t^B(S_{t+1},a^*)) -Q^*(s_t,a_t)\Big)\\
        &\mspace{36mu}+ \textbf{1}_{Y_{t+1}=0}X_t(s_t,a_t) \\
        F_t(s,a) &:= 0 \ \text{whenever } (s,a)\neq(s_t,a_t)
    \end{align*}
    with $a^*=\argmax_{a'\in\cA_{S_{t+1}}}Q^A(S_{t+1},a').$ According to [\cite{lyle19} Proposition 1] projection $\Pi_C$ is mean-preserving, i.e $\E[\Pi_C\nu]=\E[\nu]$ for when $\nu$ is a distribution supported within $[\theta_1,\theta_m]$. This is the case for every $\hat\eta_*$ as in Algorithm \ref{alg:CDRL}, which can be seen as following. Assume $\eta_t^A(s_t,a_t),\eta_t^B(s_t,a_t)\in\cF_{C,m}$. Then also 
    \[\nu= \beta^A_{t+1}\eta_t^A(S_{t+1},a^*)+(1-\beta^A_{t+1})\eta_t^B(S_{t+1},a^*))\in\cF_{C,m},\]
    and suppose $\nu=\sum_{i=1}^mp_i\delta_{\theta_i}$ for some $p_i$. Then
    \[\hat\eta_*:=b_{R_t,\gamma}\#\nu = \sum_{i=1}^mp_i\delta_{R_t+\gamma\theta_i}.\]
    But now
    \[\theta_1\leq \frac{R_{min}}{1-\gamma}\leq \frac{R_{min}}{1-\gamma}\leq \theta_m\]
    (Assumption $(ii)$) guarantees that
    \[\theta_1\leq R_t+\gamma\theta_i\leq \theta_m\quad \forall\ i\in\{1,\dots,m\}\]
    and $\hat\eta_*$ is supported within $[\theta_1,\theta_m]$.
    Similarly for a realized transition with $(R_t,S_{t+1})=(r_t,s_{t+1})$, we have for the expected value of the distribution
   \begin{align*}
        &\E\left[b_{r_t, \gamma} \#\Big(\beta^A_{t+1}\eta_t^A(s_{t+1},a^*+(1-\beta^A_{t+1})\eta_t^B(s_{t+1},a^*))\Big)\right]\\
        =&r_t + \gamma(\beta^A_{t+1}Q_t^A(s_{t+1},a^*)+(1-\beta^A_{t+1})Q_t^B(s_{t+1},a^*)).
   \end{align*}
   Hence, the expected values of the return distributions $\eta_t^A$ subtracted by $Q^*$ indeed evolve as 
   \[X_{t+1}(s,a)=(1-\alpha_t(s,a))X_t(s,a)+\alpha_t(s,a)F_t(s,a).\]
   We now proceed similarly as in \cite{hasselt2010} to show that the conditions of Lemma \ref{lem:stochastic approximation lemma} are satisfied. \\
   We first show that $\V[F_t(s,a)|\cF_t]$ is bounded for all $(s,a)\in\cS\times\cA$ and therefore satisfies $\V[F_t(s,a)|\cF_t]\leq K(1+\kappa\|X_t\|_\infty)$ as required. Since the rewards were assumed to be bounded there is an $\bar R>0$ such that $|r|,|\theta_1|,|\theta_m|\leq\bar R\ \forall r\in\cR.$ Hence, we have
   \begin{align*}
   |F_t(s_t,a_t)|\leq&|R_t + \gamma(\beta^A_{t+1}Q_t^A(S_{t+1},a^*)+(1-\beta^A_{t+1})Q_t^B(S_{t+1},a^*))-Q^*(s_t,a_t)|\\
   +&|X_t(s_t,a_t)|\\
   \leq& \bar R + 3\bar R + 2\frac{\bar R}{1-\gamma}.
   \end{align*}
   Next, we need to show that $\|\E[F_t\mid\cF_t]\|_\infty\leq \kappa \|X_t\|_\infty+c_n.$ Let us therefore decompose
 \begin{align*}
       F_t(s_t,a_t)=&\textbf{1}_{Y_{t+1}=1}\Big(F_t^Q(s_t,a_t)+\gamma(1-\beta_{t+1})(Q_t^B(S_{t+1},a^*)-Q^A_t(S_{t+1},a^*))\Big)\\
       +&\textbf{1}_{Y_{t+1}=0}\alpha_t(s_t,a_t)X_t(s_t,a_t)
 \end{align*}
   with $F_t^Q(s_t,a_t):=R_t+\gamma Q^A_t(S_{t+1},a^*)-Q^*(s_t,a_t).$ This yields
   \begin{align*}
       &|\E[\textbf{1}_{Y_{t+1}=1}F_t^Q(s_t,a_t)+\textbf{1}_{Y_{t+1}=0}\alpha_t(s_t,a_t)X_t(s_t,a_t)|\cF_t]|\\
       &=|\frac{1}{2}\E[R_t+\gamma Q^A_t(S_{t+1},a^*)]-Q^*(s_t,a_t)+\frac{1}{2}X_t(s_t,a_t)|\\
       &\leq|T^*Q^A(s_t,a_t)-T^*Q^*(s_t,a_t)|+|\frac{1}{2}X_t(s_t,a_t)|\\
       &\leq\gamma \|Q^A_t-Q^*\|_\infty+\frac{1}{2}\|X_t\|_\infty \\
       &=\underbrace{(\frac{1}{2}\gamma+\frac{1}{2})}_{<1}\|X_t\|_\infty,
   \end{align*}
   since the Bellman optimality operator is a $\gamma$-contraction. Subsequently, it only remains to show that 
   \[c_t:=|\E[\textbf{1}_{Y_{t+1}=1}\gamma(1-\beta^A_{t+1})(Q_t^B(S_{t+1},a^*)-Q^A_t(S_{t+1},a^*)|\cF_t]|\]
   goes to zero almost surely. This is immediate if we verify that
   \[X_t^{BA}(s,a):=Q^B_t(s,a)-Q^A_t(s,a)\]
   goes to zero almost surely for all $(s,a)\in\cS\times\cA$ which will be achieved by another application of Lemma \ref{lem:stochastic approximation lemma}.
   We infer that
   \begin{align*}
       &X_{n+1}^{BA}(s_n,a_n)\\
       =&X_{n}^{BA}(s_n,a_n)+\alpha_n(s_n,a_n)\biggl(\\
       &\textbf{1}_{Y_{n+1}=0}\Bigl(R_n+\gamma\bigl(\beta_{n+1}^BQ^B_n(S_{n+1},b^*)+(1-\beta^B_{n+1})Q^A_n(S_{n+1},b^*)\bigr)-Q^B_n(s_n,a_n)\Bigr) \\
       - &\textbf{1}_{Y_{n+1}=1}\Bigl(R_n+\gamma\bigl(\beta_{n+1}^AQ^A_n(S_{n+1},a^*)+(1-\beta^A_{n+1})Q^B_n(S_{n+1},a^*)\bigr)-Q^A_n(s_n,a_n)\Bigr)\\
       &\biggr)\\
       =&(1-\alpha_n(s_n,a_n))X_{n}^{BA}(s_n,a_n)+\alpha_n(s_n,a_n)\biggl(\\
       &\textbf{1}_{Y_{n+1}=0}\Bigl(R_n+\gamma\bigl(\beta_{n+1}^BQ^B_n(S_{n+1},b^*)+(1-\beta^B_{n+1})Q^A_n(S_{n+1},b^*)\bigr)\Bigr) \\
       - &\textbf{1}_{Y_{n+1}=1}\Bigl(R_n+\gamma\bigl(\beta_{n+1}^AQ^A_n(S_{n+1},a^*)+(1-\beta^A_{n+1})Q^B_n(S_{n+1},a^*)\bigr)\Bigr)\\
       +&\textbf{1}_{Y_{n+1}=1}Q^B_n(s_n,a_n)-\textbf{1}_{Y_{n+1}=0} Q^A_n(s_n,a_n)\biggr)\\
       =&(1-\alpha_n(s_n,a_n))X_{n}^{BA}(s_n,a_n)+\alpha_n(s_n,a_n)\tilde F_n(s_n,a_n),
   \end{align*}
       with 
    \begin{align*}
        \tilde F_n(s_n,a_n)=&\biggl(
       \textbf{1}_{Y_{n+1}=0}\Bigl(R_n+\gamma\bigl(\beta_{n+1}^BQ^B_n(S_{n+1},b^*)+(1-\beta^B_{n+1})Q^A_n(S_{n+1},b^*)\bigr)\Bigr) \\
       - &\textbf{1}_{Y_{n+1}=1}\Bigl(R_n+\gamma\bigl(\beta_{n+1}^AQ^A_n(S_{n+1},a^*)+(1-\beta^A_{n+1})Q^B_n(S_{n+1},a^*)\bigr)\Bigr)\\
       +&\textbf{1}_{Y_{n+1}=1}Q^B_n(s_n,a_n)-\textbf{1}_{Y_{n+1}=0} Q^A_n(s_n,a_n)\biggr).
    \end{align*}
    Now, using that $Q^B_n(s_n,a_n),Q^A_n(s_n,a_n),X_{n}^{BA}(s_n,a_n),\alpha_n(s_n,a_n)$ are $\cF_n$-measurable and $Y_{n+1}$ is independent of $\cF_n$, the conditional expectation satisfies
    \begin{align*}
        |\E[\tilde F_n(s_n,a_n)\mid\cF_n]|=&\frac{1}{2}\gamma|\E[\beta_{n+1}^BQ^B_n(S_{n+1},b^*)+(1-\beta^B_{n+1})Q^A_n(S_{n+1},b^*)\\
        -&\beta_{n+1}^AQ^A_n(S_{n+1},a^*)-(1-\beta^A_{n+1})Q^B_n(S_{n+1},a^*)|\cF_n]|\\
        +&\frac{1}{2}|Q^B_n(s_n,a_n)- Q^A_n(s_n,a_n)|\\
            \leq &\frac{1}{2} \gamma \Bigl(\left|\E[\beta^B_{n+1}(Q^B_n(S_{n+1},b^*)-Q^A_n(S_{n+1},a^*))|\cF_n]\right|\\
        +&\left|\E[(1-\beta^B_{n+1})(Q^A_n(S_{n+1},b^*)-Q^B_n(S_{n+1},a^*))|\cF_n]\right|\\
        +&\left|\E[(\beta_{n+1}^B-\beta_{n+1}^A)Q^A_n(S_{n+1},a^*)|\cF_n]\right|\\
        +&\left|\E[((1-\beta^B_{n+1})-(1-\beta^A_{n+1}))Q^B_n(S_{n+1},a^*)|\cF_n]\right|\Bigr)\\
        +&\frac{1}{2}\|X_n\|_\infty
    \end{align*}
    Now if it holds $\E[Q^B_n(S_{n+1},b^*)|\cF_n]\geq\E[Q^A_n(S_{n+1},a^*)|\cF_n]$, by definition of $a^*$ we have $Q^A_n(S_{n+1},a^*)=\max_{a\in\cA_{S_{n+1}}}Q^A_n(S_{n+1},a)\geq Q^A_n(S_{n+1},b^*)$ and therefore
    \begin{align*}
        \left|\E[Q^B_n(S_{n+1},b^*)-Q^A_n(S_{n+1},a^*)|\cF_n]\right|&=\E[Q^B_n(S_{n+1},b^*)-Q^A_n(S_{n+1},a^*)|\cF_n]\\
        &\leq\E[Q^B_n(S_{n+1},b^*)-Q^A_n(S_{n+1},b^*)|\cF_n]\leq\|X^{BA}_n\|_\infty.
    \end{align*}
    Analogously, if $\E[Q^B_n(S_{n+1},b^*)|\cF_n]<\E[Q^A_n(S_{n+1},a^*)|\cF_n]$, then we have by definition of $b^*$
    \begin{align*}
        \left|\E[Q^B_n(S_{n+1},b^*)-Q^A_n(S_{n+1},a^*)|\cF_n]\right|&=\E[Q^A_n(S_{n+1},a^*)-Q^B_n(S_{n+1},b^*)|\cF_n]\\
        &\leq\E[Q^A_n(S_{n+1},a^*)-Q^B_n(S_{n+1},a^*)|\cF_n]\leq\|X^{BA}_n\|_\infty.
    \end{align*}
    Similarly, by distinguishing cases, one shows that
    \[\left|\E[Q^A_n(S_{n+1},b^*)-Q^B_n(S_{n+1},a^*)|\cF_n]\right|\leq\|X^{BA}_n\|_\infty+\frac{1}{2}\|X^{BA}_n\|.\]
    Combining the above yields
    \begin{align*}
        \left|\E[\tilde F_n(s_n,a_n)\mid\cF_n]\right|&\leq\frac{1}{2}\gamma(\beta^B_{n+1}+(1-\beta^B_{n+1}))\|X^{BA}_n\|_\infty\\        &+\underbrace{\Big|\gamma\E[(\beta_{n+1}^B-\beta_{n+1}^A)\underbrace{Q^A_n(S_{n+1},a^*)}_{<\bar R<\infty}|\cF_n]\Big|+\Big|\gamma\E[((1-\beta^B_{n+1})-(1-\beta^A_{n+1}))\underbrace{Q^B_n(S_{n+1},a^*)}_{<\bar R<\infty}|\cF_n]\Big|}_{:=\tilde c_n  \to 0,\ \text{since } |\beta^A_n-\beta^B_n| \text{ converges to 0 for }n\to\infty\text{ due to }(iv)}.
    \end{align*}
    Hence, we invoke Lemma \ref{lem:stochastic approximation lemma} to obtain convergence of $X^{BA}_t$ and thus with another application of Lemma \ref{lem:stochastic approximation lemma}, $X_t(s,a)$ converges to zero which finally implies $Q^A_t(s,a)$ (and also $Q^B(s,a)$) converges to $Q^*(s,a)$ almost surely for every $(s,a)\in\cS\times\cA.$\\\\
    Since $\cS,\cA$ are finite, for every $\varepsilon>0$, there exists a random variable $N>0$ such that for all $t>N$, we have
    \[\max_{z\in\{A,B\}}\|Q_t^z-Q^*\|_\infty<\varepsilon\quad \text{almost surely.}\]
    \textbf{Step 2: Convergence of return distributions}\\
    Suppose the MDP has a unique optimal policy $\pi^*$. Now following \cite{rowland18}, we take $\varepsilon$ to be half the minimum action gap for the optimal action-value function $Q^*=Q^{\pi^*}$, i.e.
    \[\varepsilon=\frac{1}{2}\min_{s\in\cS}(Q^{\pi^*}(s,\pi^*(s)-\max_{a\neq}Q^{\pi^*}(s,a))\]
    which is greater than zero by assumption $(v).$
    Hence, denoting the action of the deterministic optimal policy in a certain state $s$ by $\pi^*(s)$, we get
    \[\max_aQ^A_t(s,a)=\max_aQ^B_t(s,a)=\pi^*(s)\]
    for all $t>N.$
    For some initial condition $\tilde \eta_0\in\cF_{C,m}^\cS$, let now $\tilde\eta_k$ be the iterates created by a double categorical policy evaluation algorithm for the optimal policy $\pi^*$, i.e.
    \begin{align*}
        \tilde\eta_{k+1}^A(s_k,a_k)=&(1-\textbf{1}_{Y_{k+1}=1}\alpha_k(s_k,a_k))\tilde\eta_{k}(s_k,a_k)\\
+&\textbf{1}_{Y_{k+1}=1}\alpha_k(s_k,a_k)\Pi_C\bigg(b_{R_k,\gamma}\#\Big(\beta^A_{k+1}\tilde\eta^A_k(S_{k+1},\pi^*(S_{k+1}))\\
    +&(1-\beta^A_{k+1})\tilde\eta^B_k(S_{k+1},\pi^*(S_{k+1}))\Big)\bigg)\\
        \tilde\eta_{k+1}^A(s,a)=&\tilde\eta^A_{k}(s,a)\ \text{ for }(s,a)\neq(s_k,a_k).&
    \end{align*}
    and analogously for $\tilde \eta^B.$ Note that the appearing $\Y_k,\alpha_k,\beta^A_k,\beta_k^B$ are chosen to be the same as in the control case above.
    Then $\tilde\eta^A,\tilde\eta^B$ converges almost surely to the unique fixed point $\eta^*_C$ of the projected operator $\Pi_C\cT^{\pi^*}$ with respect to $\bar\ell_2$ by Lemma \ref{lem: double categorical policy evaluation}.
    Similarly to \cite{rowland18}, we now proceed by a coupling argument. 
    Denote by $\pi_k^A,\pi^B_k$ any greedy selection rule with respect to $\eta_k^A$ and $\eta_k^B$ and $A_k=\{\pi_k^A=\pi^B_k=\pi^* \text{ for all }n\geq k\}$. Then $A_k\subseteq A_{k+1}$ and by the above explanation we have $\mathbb P(A_k)\nearrow 1$. Additionally, let $B$ be the event of probability 1 for which the (double) policy evaluation algorithm converges. Now on the event $B\cup A_k$, we have
    \[\bar \ell_2^2(\tilde\eta^A_n,\eta_C^*)\to 0\quad \text{and} \quad \bar \ell_2^2(\tilde\eta^B_n,\eta_C^*)\to 0.\]
    Then by the triangle inequality it suffices to show $\bar\ell_2(\eta^A_n,\tilde\eta^A_n)\to0$ and $\bar\ell_2(\eta^B_n,\tilde\eta^B_n)\to0$ on this event too, since then the Theorem follows by $\mathbb P(B\cup A_k)\nearrow 1$.\\
    To prove this we will again apply Lemma \ref{lem:stochastic approximation lemma}. This time with $d=2\cdot|\cS||\cA|$, where we identify 
    \[X_n := \begin{bmatrix} \ell^2_2(\eta^A_{n},\tilde\eta^A_{n}) \\ \ell^2_2(\eta^B_{n},\tilde\eta^B_{n}) \end{bmatrix}\in \R^{2|\cS||\cA|}.\]
    Additionally, we expand the filtration by $\tilde\cF_n=\sigma(\cF_n,Y_{n+1})$ and define $\tilde\alpha_n^A(s,a)=\alpha_n(s,a)\textbf{1}_{Y_{n+1}=1}$ and $\tilde\alpha_n^B(s,a)=\alpha_n(s,a)\textbf{1}_{Y_{n+1}=0}$. By Lemma \ref{lem:RM} these steps-size sequences still fulfill the Robbins-Monro conditions.\\
    Then, writing
    \begin{align*}
        \nu^A=&\beta^A_{n+1}\eta^A_n(S_{n+1},\pi^*(S_{n+1}))
    +(1-\beta^A_{n+1})\eta^B_n(S_{n+1},\pi^*(S_{n+1}))\\
    \tilde\nu^A=&\beta^A_{n+1}\tilde\eta^A_n(S_{n+1},\pi^*(S_{n+1}))
    +(1-\beta^A_{n+1})\tilde\eta^B_n(S_{n+1},\pi^*(S_{n+1}))
    \end{align*}
    for short, for $n\geq k$, on $A_k$ we have
    \begin{align*}
             &\ell^2_2(\eta^A_{n+1}(s_n,a_n),\tilde\eta^A_{n+1}(s_n,a_n))\\
             =&(1-\tilde\alpha^A_n(s_n,a_n))^2\|\eta^A_{n}(s_n,a_n)-\tilde\eta^A_{n}(s_n,a_n)\|_{\ell_2}^2\\
+&{\tilde\alpha^A_n}(s_n,a_n)^2\|\Pi_C(b_{R_n,\gamma}\#\nu^A)-\Pi_C(b_{R_n,\gamma}\#\tilde\nu^A)\|_{\ell_2}^2\\
    +&(1-\tilde\alpha^A_n(s_n,a_n))\tilde\alpha^A_n(s_n,a_n)2\langle\eta^A_{n}(s_n,a_n)-\tilde\eta^A_{n}(s_n,a_n),\Pi_C(b_{R_n,\gamma}\#\nu^A)-\Pi_C(b_{R_n,\gamma}\#\tilde\nu^A)\rangle_{\ell_2}.
    \end{align*}
    This can be rewritten in terms of Lemma \ref{lem:stochastic approximation lemma} as
    \[X_{n+1}^A(s_n,a_n)= (1-\zeta_n^A(s_n,a_n))X_{n}^A(s_n,a_n)+\zeta_n^A(s_n,a_n)F^A_n(s_n,a_n) \]
    with $ \zeta_n^A(s_n,a_n)=2\tilde\alpha^A_n(s_n,a_n)-\tilde\alpha^A_n(s_n,a_n)^2$ and
    \begin{align*}
    F^A_n(s_n,a_n)=&\frac{1}{\zeta_n^A(s_n,a_n)}({\tilde\alpha^A_n}(s_n,a_n)^2\|\Pi_C(b_{R_n,\gamma}\#\nu^A)-\Pi_C(b_{R_n,\gamma}\#\tilde\nu^A)\|_{\ell_2}^2\\
    +&(1-\tilde\alpha^A_n(s_n,a_n))\tilde\alpha^A_n(s_n,a_n)2\langle\eta^A_{n}(s_n,a_n)-\tilde\eta^A_{n}(s_n,a_n),\\
    &\Pi_C(b_{R_n,\gamma}\#\nu^A)-\Pi_C(b_{R_n,\gamma}\#\tilde\nu^A)\rangle_{\ell_2})
    \end{align*}
    and $F_n^A(s,a)=0$ if $(s,a)\neq(s_n,a_n).$ It is mentioned that $\zeta_n^A(s_n,a_n)>0.$
    Notice that,
    \begin{equation}\label{eq:zeta}
    \begin{aligned}
                \sum_{n=1}^\infty\zeta_n^A(s_n,a_n)&=\sum_{n=1}^\infty (2\tilde\alpha^A_n(s_n,a_n)-\tilde\alpha^A_n(s_n,a_n)^2)=\infty\quad a.s.\\
        \sum_{n=1}^\infty\zeta_n^A(s_n,a_n)^2&=\sum_{n=1}^\infty4 \tilde\alpha^A_n(s_n,a_n)^2 - 4 \tilde\alpha^A_n(s_n,a_n)^3 + \tilde\alpha^A_n(s_n,a_n)^2<\infty \quad a.s.
    \end{aligned}
    \end{equation}
    Finally we have
    \begin{align*}
        |F^A_n(s_n,a_n)|\leq&\frac{1}{\zeta_n^A(s_n,a_n)}({\tilde\alpha^A_n}(s_n,a_n)^2\gamma\bar\ell^2_2(\beta^A_{n+1}\eta^A_n
    +(1-\beta^A_{n+1})\eta^B_n,\beta^A_{n+1}\tilde\eta^A_n
    +(1-\beta^A_{n+1})\tilde\eta^B_n)\\
    +&(1-\tilde\alpha^A_n(s_n,a_n))\tilde\alpha^A_n(s_n,a_n)2\sqrt{\gamma}|\langle\eta^A_{n}-\tilde\eta^A_{n},\\
    &\beta^A_n\eta^A_n
    +(1-\beta^A_n)\eta^B_n-\beta^A_n\tilde\eta^A_n
    -(1-\beta^A_n)\tilde\eta^B_n\rangle_{\bar\ell_2}|)\\
    \leq&\frac{1}{\zeta_n^A(s_n,a_n)}({\tilde\alpha^A_n}(s_n,a_n)^2\gamma\max_{z\in\{A,B\}}\bar\ell^2_2(\eta^z_{n},\tilde\eta^z_{n})\\
    +&(1-\tilde\alpha^A_n(s_n,a_n))\tilde\alpha^A_n(s_n,a_n)2\sqrt{\gamma}\max_{z\in\{A,B\}}\bar\ell^2_2(\eta^z_{n},\tilde\eta^z_{n}))\\
    =&\frac{{\tilde\alpha^A_n}(s_n,a_n)^2\gamma
    +(1-\tilde\alpha^A_n(s_n,a_n))\tilde\alpha^A_n(s_n,a_n)2\sqrt{\gamma}}{2\tilde\alpha^A_n(s_n,a_n)-\tilde\alpha^A_n(s_n,a_n)^2}\max_{z\in\{A,B\}}\bar\ell^2_2(\eta^z_{n},\tilde\eta^z_{n})\\
    \leq&\frac{
    (2\tilde\alpha^A_n(s_n,a_n)-\tilde\alpha^A_n(s_n,a_n)^2)\sqrt{\gamma}}{2\tilde\alpha^A_n(s_n,a_n)-\tilde\alpha^A_n(s_n,a_n)^2}\max_{z\in\{A,B\}}\bar\ell^2_2(\eta^z_{n},\tilde\eta^z_{n})\\
  \leq& \sqrt{\gamma}\max_{z\in\{A,B\}}\bar\ell^2_2(\eta^z_{n},\tilde\eta^z_{n})=\sqrt{\gamma}\|X_n\|_\infty
    \end{align*}
    where we used regularity and 1/2-homogeneity of the $\ell_2$ metric as described in [\cite{drl} Section 4.6] as well as that $\Pi_C$ is a non-expansion in $\ell_2$ and
   \begin{align*}
        &|\langle u,\beta u +(1-\beta) v\rangle|=\beta\langle u,u\rangle+(1-\beta)|\langle u,v\rangle|\leq\beta \max(\|u\|^2,\|v\|^2)+(1-\beta)\|u\|\|v\|\\
        \leq&\max(\|u\|^2,\|v\|^2) 
   \end{align*}
    by the Cauchy-Schwarz inequality. Further, by the above the Variance also fulfills
    \begin{align*}
        \V[F^A_n(s_n,a_n)|\tilde\cF_n]=\E[F^A_n(s_n,a_n)^2|\cF_n]-\E[F^A_n(s_n,a_n)|\tilde\cF_n]^2\\
        \leq 2 (\sqrt{\gamma}\max_{z\in\{A,B\}}\bar\ell^2_2(\eta^z_{n},\tilde\eta^z_{n}))^2\\
        \leq2 \gamma \sup_{\eta,\eta\in\cF_{C,m}^\cS}\bar\ell^4_2(\eta,\eta')<\infty.
    \end{align*}
    Therefore, by Lemma \ref{lem:stochastic approximation lemma} we obtain convergence $\bar\ell_2(\eta^A_n,\tilde\eta^A_n)\to0$ and $\bar\ell_2(\eta^B_n,\tilde\eta^B_n)\to0$ on $A_k$. As already described above, this results in 
    \[\bar\ell_2(\eta^A_n,\eta_C^*)\to0\quad  \text{and}\quad \bar\ell_2(\eta^B_n,\eta_C^*)\to0\quad \text{almost surely.}\]
\end{proof}

\begin{proof}[Proof of Lemma \ref{lem: double categorical policy evaluation}]
Let the filtration be given by $\cF_t=\sigma(\eta^A_0,\eta^B_0,s_0,a_0,\alpha_0,R_0,S_1,Y_1,\beta^A_1,\beta^B_1\dots,s_t,a_t,\alpha_t,Y_{t+1}),$ where $(Y_{n})_{n\in\N}$ is an iid sequence of \textit{Bernoulli(0.5)} random variables, independent of all other appearing random variables, such that A is updated when $Y_{n+1}=1$.
To clarify, abbreviating
\begin{align*}
    \nu^A&=\beta_{t+1}^A\eta^A_{t}(S_{t+1},A_{t+1})+(1-\beta_{t+1}^A)\eta^B_{t}(S_{t+1},A_{t+1})\\
    \nu^B&=\beta_{t+1}^B\eta^B_{t}(S_{t+1},A_{t+1})+(1-\beta_{t+1}^B)\eta^A_{t}(S_{t+1},A_{t+1})\quad \text{where}\\
    A_{t+1}&\sim\pi(\cdot;S_{t+1}),
\end{align*}
we are confronted with the updates
    \begin{align*}
        \eta^A_{t+1}(s,a)&=\eta^A_{t}(s,a)+\alpha_t(s,a)\textbf{1}_{Y_{t+1}=1}(\Pi_C(b_{R_t,\gamma}\#\nu^A)-\eta^A_{t}(s,a))\\
        \eta^B_{t+1}(s,a)&=\eta^B_{t}(s,a)+\alpha_t(s,a)\textbf{1}_{Y_{t+1}=0}(\Pi_C(b_{R_t,\gamma}\#\nu^B)-\eta^B_{t}(s,a)).
    \end{align*}
 As in the proof above, define $\tilde\alpha_n^A(s,a)=\alpha_n(s,a)\textbf{1}_{Y_{n+1}=1}$ and $\tilde\alpha_n^B(s,a)=\alpha_n(s,a)\textbf{1}_{Y_{n+1}=0}$. By Lemma \ref{lem:RM} these steps-size sequences still fulfill the Robbins-Monro conditions. Also note that as in step 2 of the proof of Theorem \ref{thm:double categorical control}, $Y_{t+1}$ is $\cF_{t}$-measurable and hence so is $\tilde\alpha_t^{A/B}$.
In order to align this with Lemma \ref{lem:stochastic approximation lemma}, we rewrite
\begin{align*}
    &X^A_{t+1}(s,a)=\ell_2^2(\eta^A_{t+1}(s,a),\eta_C(s,a))\\
    =&(1-\tilde\alpha^A_t(s,a))^2\|\eta^A_{t}(s,a)-\eta_C(s,a)\|_{\ell_2}^2\\
+&{\tilde\alpha^A_t}(s,a)^2\|\Pi_C(b_{R_t,\gamma}\#\nu^A)-\eta_C(s,a)\|_{\ell_2}^2\\
    +&(1-\tilde\alpha^A_t(s,a))\tilde\alpha^A_t(s,a)2\langle\eta^A_{t}(s,a)-\eta_C(s,a),\Pi_C(b_{R_t,\gamma}\#\nu^A)-\eta_C(s,a)\rangle_{\ell_2}\\
    =&(1-\zeta_t^A(s,a))X_{t}^A(s,a)+\zeta_t^A(s,a)F^A_t(s,a) 
\end{align*}
    with $ \zeta_t^A(s,a)=2\tilde\alpha^A_t(s,a)-\tilde\alpha^A_t(s,a)^2$,
    \[X_t := \begin{bmatrix} \ell^2_2(\eta^A_{t},\eta_C) \\ \ell^2_2(\eta^B_{t},\eta_C) \end{bmatrix}\in \R^{2|\cS||\cA|}\]
    and
    \begin{align*}
    F^A_t(s,a)=&\frac{1}{\zeta_t^A(s,a)}\textbf{1}_{\tilde\alpha^A_t(s,a)>0}({\tilde\alpha^A_t}(s,a)^2\ell_2^2(\Pi_C(b_{R_t,\gamma}\#\nu^A),\eta_C(s,a))\\
    +&(1-\tilde\alpha^A_t(s,a))\tilde\alpha^A_t(s,a)2\langle\eta^A_{t}(s,a)-\eta_C(s,a),\Pi_C(b_{R_t,\gamma}\#\nu^A)-\eta_C(s,a)\rangle_{\ell_2}).
    \end{align*}
    As in Equation (\ref{eq:zeta}), the sequence $ \zeta_t^A(s,a)$ fulfills the Robbins-Monro condition. Additionally, note that there exists $K>0$, such that $\ell_2^2(\Pi_C(b_{R_t,\gamma}\#\nu^A),\eta_C(s,a))<K$ independent of $s,a,t$. Further, observe that
    \[c_t:=\max_{z\in\{A,B\}} \frac{1}{\zeta_t^z(s,a)}\textbf{1}_{\tilde\alpha^z_t(s,a)>0}{\tilde\alpha^z_t}(s,a)^2 K\to 0\ \text{for } t\to\infty\text{ almost surely}.\]
We use that $\Pi_C$ is mean-preserving [\cite{lyle19} Proposition 1] for discrete distributions supported within $[\theta_1,\theta_m]$, which is satisfied by $b_{R_t,\gamma}\#\nu^A$, due to Assumption $(ii)$ and $\nu^A\in\cF_m$. Together with the fact that $\Pi_C\cT^\pi$ is a $\sqrt{\gamma}$-contraction with respect to $\bar\ell_2$  and the Cauchy-Schwarz inequality, we have
\begin{align*}
    &|\E[\langle\eta^A_{t}(s,a)-\eta_C(s,a),\Pi_C(b_{R_t,\gamma}\#\nu^A)-\eta_C(s,a)\rangle_{\ell_2}|\cF_t]|\\
    =&|\langle\eta^A_{t}(s,a)-\eta_C(s,a),\E[\Pi_C(b_{R_t,\gamma}\#\nu^A)|\cF_t]-\eta_C(s,a)\rangle_{\ell_2}|\\
    =&|\langle\eta^A_{t}(s,a)-\eta_C(s,a),\E[b_{R_t,\gamma}\#\nu^A)|\cF_t]-\eta_C(s,a)\rangle_{\ell_2}|\\
    =&|\langle\eta^A_{t}(s,a)-\eta_C(s,a),\Pi_C\cT^\pi(\beta_{t+1}^A\eta^A_{t}+(1-\beta_{t+1}^A)\eta^B_{t})(s,a)-(\Pi_C\cT^\pi\eta_C)(s,a)\rangle_{\ell_2}|\\
    \leq&\sqrt{\gamma}|\langle\eta^A_{t}-\eta_C,(\beta_{t+1}^A\eta^A_{t}+(1-\beta_{t+1}^A)\eta^B_{t})-\eta_C\rangle_{\bar\ell_2}|\\
    \leq&\sqrt{\gamma}(\beta_{t+1}^A\bar\ell_2^2(\eta^A_{t},\eta_C)+(1-\beta_{t+1}^A)|\langle\eta^A_{t}-\eta_C,\eta^B_{t}-\eta_C\rangle_{\bar\ell_2}|)\\
    \leq&\sqrt{\gamma}(\beta_{t+1}^A\max_{z\in\{A,B\}}\bar\ell_2^2(\eta^z_{t},\eta_C)+(1-\beta_{t+1}^A)\|\eta^A_{t}-\eta_C\|_{\bar\ell_2}\|\eta^B_{t}-\eta_C\|_{\bar\ell_2}|)\\
    \leq&\sqrt{\gamma}\max_{z\in\{A,B\}}\bar\ell_2^2(\eta^z_{t},\eta_C)\\
    =&\sqrt{\gamma}\|X_t\|_\infty.
\end{align*}
    In total, this yields
    \begin{align*}
        &|\E[F^A_t(s,a)|\cF_t]|\\
        &\leq \frac{1}{\zeta_t^A(s,a)}\textbf{1}_{\tilde\alpha^A_t(s,a)>0}{\tilde\alpha^A_t}(s,a)^2 K + \frac{1}{\zeta_t^A(s,a)}\textbf{1}_{\tilde\alpha^A_t(s,a)>0}(1-\tilde\alpha^A_t(s,a))\tilde\alpha^A_t(s,a)2\sqrt{\gamma}\|X_t\|_\infty.\\
        &\leq c_t + \sqrt{\gamma}\|X_t\|_\infty.
    \end{align*}    
    Since $\bar\ell_2(\eta,\eta')<K$ for every $\eta,\eta'\in\cF_{C,m}^{\cS\times\cA}$ some $K>0$, the conditional variance $\V[F^A_t|\cF_t]$ can be bounded uniformly in $t$. \\
    Therefore, the requirements of Lemma \ref{lem:stochastic approximation lemma} are fulfilled, and its application yields $X_t^A(s,a)=\ell_2^2(\eta_t^A(s,a),\eta_C(s,a))\to0$ and $X_t^B(s,a)=\ell_2^2(\eta_t^B(s,a),\eta_C(s,a))\to0$. Hence, also $\eta_t^A,\eta_t^B$ converge to $\eta_C$ with respect to $\bar\ell_2$.
\end{proof}




\end{document}
